===== BEGIN FILE: 2604.25833v2.tex =====
\documentclass[
 twocolumn,
 superscriptaddress,
 preprintnumbers,
 nofootinbib,
 amsmath,amssymb,
 aps,
 prl
 ]{revtex4-2}
\usepackage{slashed,bm}
\usepackage{url}

\usepackage[colorlinks = true,
            linkcolor = blue,
            urlcolor  = blue,
            citecolor = blue,
            anchorcolor = blue]{hyperref}

















\newcommand{\sign}{\mathop{\mathrm{sign}}\nolimits}
\newcommand{\tr}{\mathop{\mathrm{tr}}\nolimits}
\newcommand{\NP}{\slashed{\mathcal{P}}}
\newcommand{\BS}{\mathbb{S}}
\newcommand\HS{\mathrm{S}}
\newcommand{\HP}{\mathrm{H}}
\newcommand{\HPL}{\mathrm{H}}
\def\deltamx{\delta(1-x)}
\def\ggp{p_{gg}(x)}
\def\ggm{p_{gg}(-x)}
\def\ca{C_A}
\def\cf{C_F}
\def\nf{n_f}
\def\D#1{D_#1}
\def\z#1{\zeta_#1}
\def\frct#1#2{\mbox{\small{$\displaystyle\frac{#1}{#2}$}}}
\def\frkt#1#2{\mbox{\large{$\frac{#1}{#2}$}}}




























































\allowdisplaybreaks

\begin{document}

\preprint{PUBDB-2026-01395, LTH 1415}

\title{Four-Loop Gluon Anomalous Dimension of General Lorentz Spin: Transcendental Part}

\author{B.A.~Kniehl}
\email{kniehl@desy.de}
\affiliation{II. Institute for Theoretical Physics, Hamburg University,
 D-22761 Hamburg, Germany}
\author{S.-O.~Moch}
\email{sven-olaf.moch@desy.de}
\affiliation{II. Institute for Theoretical Physics, Hamburg University,
 D-22761 Hamburg, Germany}
\author{V.N.~Velizhanin}
\email{vitaly.velizhanin@desy.de}
\affiliation{II. Institute for Theoretical Physics, Hamburg University,
 D-22761 Hamburg, Germany}
\author{A.~Vogt}
\email{andreas.vogt@liverpool.ac.uk}
\affiliation{II. Institute for Theoretical Physics, Hamburg University,
 D-22761 Hamburg, Germany}
\affiliation{Department of Mathematical Sciences, University of Liverpool,
 Liverpool L69 3BX, United Kingdom}




\begin{abstract}
  We consider the anomalous dimension $\gamma_{gg}^{(3)}(N)$ of the twist-two gluon operator of arbitrary Lorentz spin $N$ in the quark flavor singlet sector of a general gauge theory at four loops and construct its contribution proportional to $\zeta(3)$ in analytic form by applying the Lenstra-Lenstra-Lov\'{a}sz algorithm to the available low-$N$ moments.
We exploit generalized Gribov--Liptov reciprocity, establish new self-tuning relations for the anomalous dimension matrix of the singlet sector, and inject information from $\mathcal{N}=1,4$ supersymmetric Yang--Mills theories.
We also present the contribution to the rational part of $\gamma_{gg}^{(3)}(N)$ with color factor $C_F^2n_f^2$.
Exact contributions to the four-loop splitting function $P_{gg}^{(3)}(x)$ hence resulting via inverse Mellin transformation help us to reduce theoretical uncertainties in scaling violations of parton distribution functions in QCD.
\end{abstract}


\pacs{}


\maketitle




The exploitation of the physics potentials of future high-precision measurements at the CERN Large Hadron Collider, the BNL Electron-Ion Collider, and other hadron-beam facilities critically depends on our ability to deepen our understanding of the proton's parton density functions (PDFs), which enter every cross section prediction in the parton model \cite{Bjorken:1969ja} of quantum chromodynamics (QCD) \cite{Fritzsch:1973pi}.
While PDFs are genuinely nonperturbative objects, their scaling violations, encoded in Dokshitzer--Gribov--Lipatov--Altarelli--Parisi (DGLAP) \cite{Gribov:1972ri,Gribov:1972rt,Altarelli:1977zs,Dokshitzer:1977sg} evolution, are amenable to perturbation theory in the strong-coupling constant $\alpha_s$.
This and PDF universality, both consequences of the factorization theorem \cite{Collins:1989gx}, provide the theoretical foundation both for determinations of PDFs from global fits to experimental data and their predictive power for hadron-induced processes.
At the same time, this allows for high-precision determinations of $\alpha_s$ \cite{Alekhin:2025qdj} and heavy-quark masses \cite{Alekhin:2010sv,Alekhin:2012un}.

The DGLAP evolution kernels consist of splitting functions, $P_{ij}(x)$, which, in a way, measure the probability of parton (gluon or quasi-massless quark) $j$ to inclusively branch into collinear parton $i$ carrying fraction $x$ of $j$'s momentum.
The perturbative expansion
\begin{equation}
P_{ij}(x)=\sum_{n=0}^{\infty}a_s^{n+1}P_{ij}^{(n)}(x)\,,
\end{equation}
with $a_s=\alpha_s(\mu)/(4\pi)$, is a topic of old vintage, dating back to the early 1970s, when the leading order (LO) was explored \cite{Gribov:1972ri,Gribov:1972rt,Altarelli:1977zs,Dokshitzer:1977sg,Gross:1973ju,Gross:1974cs}, and has enjoyed top priority on the agenda of the world-wide particle physics community ever since.
Full analytic expressions are available at next-to-LO (NLO) \cite{Floratos:1977au,GonzalezArroyo:1979df,Gonzalez-Arroyo:1979kjx,Curci:1980uw,Furmanski:1980cm} and N${}^2$LO \cite{Larin:1993vu,Larin:1996wd,Moch:2004pa,Vogt:2004mw,Ablinger:2014nga,Blumlein:2021enk,Blumlein:2022gpp,Gehrmann:2023ksf}, while only partial results are yet known at N${}^3$LO \cite{Gracey:1994nn,Velizhanin:2011es,Velizhanin:2014fua,Davies:2016jie,Moch:2017uml,Davies:2017hyl,Moch:2018wjh,Das:2020adl,Moch:2021qrk,Falcioni:2023luc,Gehrmann:2023cqm,Falcioni:2023vqq,Falcioni:2023tzp,Moch:2023tdj,Gehrmann:2023iah,Falcioni:2024xyt,Falcioni:2024xav,Falcioni:2024qpd,Kniehl:2025jfs,Kniehl:2025ttz,Falcioni:2025hfz,Kniehl:2026axe,Gehrmann:2026qbl}.
Modern evaluations exploit the fact that the Mellin transforms
\begin{equation}
  \gamma_{ij}(N)=-\int_0^1 dx\ x^{N-1} P_{ij}(x)\,,
\label{eq:mel}
\end{equation}
are anomalous dimensions of local composite operators of twist two and Lorentz spin $N$, constructed from quark fields, gluon field strength tensors, and a definite number of covariant derivatives, namely $N-1$ for quark operators and $N-2$ for gluon operators \cite{Dokshitzer:1977sg}.
Specifically, one renormalizes off-shell matrix elements of such operators in the $\overline{\mathrm{MS}}$ scheme.
At N$^3$LO, one must deal with four-loop Feynman diagrams, whose complexity rapidly grows with increasing $N$ until the limits of available computer algebra systems and computing resources are reached.
In turn, from finite sets of Mellin moments, only approximate results for $P_{ij}$ can be gained \cite{Larin:1996wd,Moch:2017uml,Moch:2023tdj,Falcioni:2024xyt,Falcioni:2025hfz}, which generally become insufficient in the high-energy limit, at very low $x$ values (see, {\it e.g.}, Figs.~5 and 6 in Ref.~\cite{Moch:2017uml}), as probed by the LHC.
This strongly motivates us to recover all-$N$ results for $\gamma_{ij}^{(3)}$ as much as possible.

Given the PDFs of $f$-flavored quarks and antiquarks and the gluon, $q_f$, $\bar{q}_f$, and $g$, the quark sector is organized in terms of flavor asymmetries $q_{\mathrm{ns},ff^\prime}^{\pm}=(q_f\pm \bar{q}_f)-(q_{f^\prime}\pm \bar{q}_{f^\prime})$, flavor non-singlet valence $q_{\mathrm{ns}}^{\mathrm{v}}=\sum_{f=1}^{n_f}(q_f-\bar{q}_f)$, and flavor singlet $q_{\mathrm{s}}=\sum_{f=1}^{n_f}(q_f+\bar{q}_f)$, where $n_f$ is the number of active quark flavors.
Then the DGLAP evolution decouples for $q_{\mathrm{ns},ff^\prime}^{\pm}$ and $q_{\mathrm{ns}}^{\mathrm{v}}$, leaving a $2\times2$ matrix equation for $q_{\mathrm{s}}$ and $g$.
This involves seven distinct splitting functions, $P_{\mathrm{ns}}^{\pm}$, $P_{\mathrm{ns}}^{\mathrm{v}}$, $P_{qq}$, $P_{qg}$, $P_{gq}$, $P_{gg}$.
Furthermore, one decomposes $P_{\mathrm{ns}}^{\mathrm{v}}=P_{\mathrm{ns}}^-+P_{\mathrm{s}}$ and $P_{qq}=P_{\mathrm{ns}}^++P_{\mathrm{ps}}$, involving quark sea and pure singlet.
This very structure is mapped to Mellin space via Eq.~\eqref{eq:mel}.

In this letter, we focus attention on $\gamma_{gg}^{(3)}$, at N$^3$LO.
This is known to break down in rational and transcendental parts as \cite{Moch:2021qrk,Moch:2023tdj,Falcioni:2024qpd,Falcioni:2025hfz}
\begin{equation}
\gamma_{gg}^{(3)}=
\gamma_{gg,\mathrm{rat}}^{(3)}
+\zeta_3\gamma_{gg,\zeta_3}^{(3)}
+\zeta_4\gamma_{gg,\zeta_4}^{(3)}
+\zeta_5\gamma_{gg,\zeta_5}^{(3)}\,,
\label{eq:zeta3}
\end{equation}
where $\zeta_k=\zeta(k)$ is Riemann's zeta function.
A term proportional to $\zeta_2=\pi^2/6$ is prohibited by the no-$\pi^2$ theorem \cite{Davies:2017hyl,Jamin:2017mul,Baikov:2018wgs} of the $\overline{\mathrm{MS}}$ scheme.
Presently, $\gamma_{gg,\zeta_5}^{(3)}$ \cite{Moch:2018wjh,Falcioni:2024qpd}, $\gamma_{gg,\zeta_4}^{(3)}$ \cite{Davies:2017hyl}, the $n_f^3$ \cite{Davies:2016jie} and $n_f^2$ \cite{Falcioni:2024qpd} terms of $\gamma_{gg,\zeta_3}^{(3)}$, and the $n_f^3$ term of $\gamma_{gg,\mathrm{rat}}^{(3)}$ \cite{Davies:2016jie} are known for all values of $N$, while the residual terms, namely the $n_f^{1,0}$ terms of $\gamma_{gg,\zeta_3}^{(3)}$ and the $n_f^{2,1,0}$ terms $\gamma_{gg,\mathrm{rat}}^{(3)}$, are only available for $N=2,4,\ldots,22$ \cite{Moch:2021qrk
,Moch:2023tdj
,Falcioni:2024qpd
,Falcioni:2025hfz
}.
The results for odd values of $N$ are inaccessible via the operator product expansion.
In the following, we analytically reconstruct the all-$N$ result for $\gamma_{gg,\zeta_3}^{(3)}$ utilizing advanced techniques of number theory in combination with an understanding of the appearing classes of special functions, conjectured theorems, and educated guesses.
This completes our knowledge of the transcendental part of $\gamma_{gg}^{(3)}$ for arbitrary value of $N$.
Recently, the same has been done for $\gamma_{\mathrm{ns}}^{(3)\pm}$, $\gamma_{\mathrm{s}}^{(3)}$ \cite{Kniehl:2025jfs,Kniehl:2026axe,Gehrmann:2026qbl}, $\gamma_{\mathrm{ps}}^{(3)}$ \cite{Falcioni:2025hfz}, and partly also for $\gamma_{qg}^{(3)}$ and $\gamma_{gq}^{(3)}$, up to $n_f^1$ and $n_f^0$ terms, respectively \cite{Falcioni:2025hfz}.

Detailed inspection of the known all-$N$ expressions for $\gamma_{ij}^{(n)}$, at $n+1$ loops, and parts thereof reveals that the special functions appearing therein are exhausted by the nested harmonic sums \cite{Vermaseren:1998uu},
\begin{equation}
S_{m_1,\ldots,m_M}(N)=\sum^{N}_{i=1} \frct{[\sign(m_1)]^{i}}{i^{\vert m_1\vert}}
\,S_{m_2,\ldots,m_M}(i)\,,
\label{eq:hs}
\end{equation}
where $S(N)=1$ and $m_j\in\mathbb{Z}\backslash\{-1,0\}$, with weight (transcendentality) $w=\sum_{j=1}^M|m_j|\le2n+1$.
At weight $w$, there are $[(1-\sqrt{2})^w+(1+\sqrt{2})^w]/2$ of them.
In addition, their counterparts with shifted arguments $N\pm1$ and $N\pm2$ may occur.
These may be rewritten as linear combinations of $D_1$, $D_2$, $\eta$, and $\nu$, having $w=1$, their powers, and products of these terms with nested harmonic sums of argument $N$, where
\begin{equation}
\eta=D_0-D_1\,,\qquad \nu=D_{-1}-D_2\,,\qquad D_k=\frac{1}{N+k}\,.
\label{eq:etanu}
\end{equation}
Notice that $\zeta_k$ has $w=k$ and its appearance as an overall factor correspondingly reduces the maximum weight of the function basis used for the ansatz of analytic reconstruction.
Fortunately, the coefficients in the ansatz empirically come as relatively harmless fractions, made of integers of moderate size modulo some powers of 2 and 3.
Therefore, we may attempt to determine them using the Lenstra-Lenstra-Lov\'{a}sz (LLL) algorithm \cite{Lenstra82factoringpolynomials} as implemented in the program package \texttt{fplll} \cite{fplll}, which allows us to solve systems of
linear equations with far more unknowns than equations.
A reliable success indicator is that several coefficients vanish and that the remaining ones are simple fractions.
A typical example of how this comes about is presented in the Supplemental Material~\ref{supplement:1}.
This method for reconstructing all-$N$ expressions of anomalous dimensions from a finite set of Mellin moments was pioneered in $\mathcal{N}=4$ supersymmetric Yang--Mills theory (SYM) \cite{Velizhanin:2010cm} and subsequently applied in ${\mathcal N}=4$ \cite{Velizhanin:2013vla,Marboe:2014sya,Marboe:2016igj,Kniehl:2020rip,Kniehl:2021ysp,Velizhanin:2021bdh,Kniehl:2023bbk,Kniehl:2024tvd} and $\mathcal{N}=2$ SYM \cite{Kniehl:2023bbk} and also in QCD \cite{Velizhanin:2012nm,Davies:2016jie,Moch:2017uml,Kniehl:2025jfs,Kniehl:2025ttz,Falcioni:2025hfz,Kniehl:2026axe}.
Prior to applying the LLL algorithm \cite{Lenstra82factoringpolynomials,fplll}, we may fix the coefficients of $\eta^{4,3,2}$ and $\nu^4$ using the Chinese remainder theorem \cite{china} as originally done in Ref.~\cite{Moch:2014sna}.
According to the na\"{\i}ve weight counting rules mentioned above, the coefficient of $\zeta(3)$ has $w=4$, so that there are hundreds of nested harmonic sums, which is far beyond what the LLL algorithm can manage with just eleven $\gamma_{gg,\zeta_3}^{(3)}$ inputs.

This situation would be hopeless, if it were not for Gribov--Lipatov reciprocity, originally observed at LO as $P_{ii}^{(0)}(x)=-xP_{ii}^{(0)}(1/x)$ or, equivalently, as the quasi-invariance of $\gamma_{ii}^{(0)}(N)$ under the mapping $N\to-1-N$ \cite{Gribov:1972rt}.
This incorporates a similar relation between time-like and space-like structure functions previously established by Drell, Levy, and Yan \cite{Drell:1969jm,Blumlein:2000wh}.
In the non-singlet sector, the $N\leftrightarrow-1-N$ symmetry was recovered also beyond LO, albeit in the alleviated sense that it no longer holds for the full anomalous dimension $\gamma(N)$, but only for the so-called reciprocity-respecting part $\mathcal{P}(N)$ of its representation in terms of nested harmonic sums in the large-$N$ limit, which is known to satisfy the self-tuning relation \cite{Dokshitzer:2005bf,Dokshitzer:2006nm,Basso:2006nk},
\begin{equation}
\gamma(N)=\mathcal{P}(N-\gamma(N)-\beta(a_s)/a_s)\,,
\label{eq:self}
\end{equation}
where
\begin{equation}
\mu^2\frct{\mathrm{d}a_s}{\mathrm{d}\mu^2}=
\beta(a_s)=-\sum_{n=0}^\infty b_n\,a_s^{n+2}\,,
\end{equation}
is the $\beta$ function of QCD \cite{Gross:1973id,Politzer:1973fx,Jones:1974mm,Caswell:1974gg,Tarasov:1980au,Larin:1993tp,vanRitbergen:1997va,Chetyrkin:2004mf,Czakon:2004bu,Luthe:2016ima,Baikov:2016tgj,Herzog:2017ohr,Luthe:2017ttg}.
This entails the decomposition $\mathcal{P}(N)=\gamma(N)+\NP(N)$, where the coefficients $\NP^{(n)}(N)$ of the non-reciprocity-respecting remainder,
\begin{equation}
  \NP(N)=\sum_{n=1}^\infty a_s^{n+1}\NP^{(n)}(N)\,,
\end{equation}
are expressed entirely in terms of $\gamma^{(k)}$, its derivatives with respect to $N$, and $b_k$ with $k<n$.
Explicit expressions for $\NP^{(n)}$ with $n=1,2,3$ may be found in Refs.~\cite{Kniehl:2025jfs,Kniehl:2026axe}.
In other words, $\NP^{(3)}$ is fully determined by known results from N${}^2$LO and below.
The salient point is that the function space thus collapses to the much smaller reciprocity-respecting subspace of binomial harmonic sums \cite{Vermaseren:1998uu},
\begin{equation}
\BS_{m_1,\ldots,m_M}\!(N)\!=\!(-1)^N\!\sum_{i=1}^{N}\!(-1)^i\!\binom{N}{i}\!\!\binom{N+i}{i}\!S_{m_1,...,m_M}(i)\,,
\end{equation}
where $m_j\in\mathbb{N}$.
In fact, there are only $2^{w-1}$ such sums with weight $w$.
Similarly to the case of nested harmonic sums, we have to allow for the binomial harmonic sums to be dressed by powers of $\eta$ and $\nu$, which are both reciprocity respecting.

In the singlet sector, Eq.~\eqref{eq:self} does not hold for the individual entries of the matrix
\begin{equation}
\hat{\gamma}=\left( \begin{array}{cc}
         \gamma_{qq} & \gamma_{qg} \\
         \gamma_{gq} & \gamma_{gg} \end{array} \right)\,,
\label{eq:matrix}
\end{equation}
where $\gamma_{qq}=\gamma_{\mathrm{ns}}^++\gamma_{\mathrm{ps}}$.
However, detailed inspection reveals that the eigenvalues, $\gamma_\pm=\tr\hat{\gamma}/2\pm[(\tr\hat{\gamma})^2/4-\det\hat{\gamma}]^{1/2}$, do satisfy Eq.~\eqref{eq:self}.
This follows from Ref.~\cite{Basso:2006nk} and further studies of large-spin systematics~\cite{Alday:2015eya}.
A similar observation was made at N${}^2$LO in Ref.~\cite{Chen:2020uvt}, mixing spacelike and timelike quantities, however.
Our formulation has the advantage of being purely spacelike.
More details will be presented in Ref.~\cite{Moch-Shumilov-to-appear}.
Because of the linearity of the decomposition into reciprocity-respecting and non-reciprocity-respecting parts, also $\gamma_++\gamma_-=\tr\hat{\gamma}=\gamma_{qq}+\gamma_{gg}$ satisfies Eq.~\eqref{eq:self}.
Exploiting the knowledge of $\gamma_{\mathrm{ns},\zeta_3}^{(3)+}$ \cite{Kniehl:2025jfs,Kniehl:2026axe,Gehrmann:2026qbl} and $\gamma_{\mathrm{ps},\zeta_3}^{(3)}$ \cite{Falcioni:2025hfz} for all $N$, we may thus analytically reconstruct $\tr\hat{\gamma}_{\zeta_3}^{(3)}$ from the known Mellin moments of $\gamma_{gg,\zeta_3}^{(3)}$ \cite{Moch:2021qrk,Moch:2023tdj,Falcioni:2024qpd,Falcioni:2025hfz} and so obtain $\gamma_{gg,\zeta_3}^{(3)}=\tr\hat{\gamma}_{\zeta_3}^{(3)}-\gamma_{\mathrm{ns},\zeta_3}^{(3)+}-\gamma_{\mathrm{ps},\zeta_3}^{(3)}$ for all $N$.

The sought-after $n_f^1$ and $n_f^0$ terms of $\gamma_{gg,\zeta_3}^{(3)}$ come with color factors $n_fC_F^3$, $n_fC_AC_F^2$, $n_fC_A^2C_F$, $n_fC_A^3$, $n_fd_{RA}^{(4)}/n_a$, $C_A^4$, and $d_{AA}^{(4)}/n_a$ \cite{Moch:2021qrk,Moch:2023tdj,Falcioni:2024qpd,Falcioni:2025hfz}, where $C_F=(n_c^2-1)/(2n_c)$, $C_A=n_c$, $d_{RA}^{(4)}=d_F^{abcd}d_A^{abcd}/n_a=n_c(n_c^2+6)/48$ and $d_{AA}^{(4)}/n_a=d_A^{abcd}d_A^{abcd}/n_a=n_c^2(n_c^2+36)/24$ with $n_a=n_c^2-1$ for color gauge group SU($n_c$).
By the conjectured maximal-transcendentality principle \cite{Kotikov:2002ab,Kotikov:2006ts}, the $w=4$ terms in the coefficients of $C_A^4$ and $d_{AA}^{(4)}/n_a$ may be gleaned respectively from the planar \cite{Beisert:2006ez,Bajnok:2008qj} and nonplanar \cite{Kniehl:2020rip,Kniehl:2021ysp,Kniehl:2024tvd} universal anomalous dimensions of the twist-two operators with general Lorentz spin $N$ at N${}^3$LO in $\mathcal{N}=4$ SYM, which serves as welcome input.
Since the quartic color factors $n_fd_{RA}^{(4)}/n_a$ and $d_{AA}^{(4)}/n_a$ appear for the first time at N${}^3$LO, their coefficients are reciprocity respecting and may be reconstructed directly.
In terms of binomial harmonic sums, we find
\begin{widetext}
\begin{eqnarray}
&&  \gamma_{gg,\zeta_3}^{(3)}\Big|_{\mathrm{quar}}=





\frct{128}{3}\*\frct{d_{RA}^{(4)}}{n_a}\*n_f\*\Big[
\frct{103}{15}
+\frct{2908}{45}\*\nu
-291\*\eta
-387\*\eta^2
+\BS_{1}\*\Big(\frct{52}{3}\*\nu +5 + 14\*\eta - 46\*\eta ^2 - 84\*\eta ^3\Big)
-\BS_{2}\*\Big(\frct{514}{15}+\frct{618}{5}\*\eta
- 4\*\nu
\nonumber\\
&&{}
- 132\*\eta ^2\Big)
+\BS_{1,1}\*\big(9\*\eta - 16\*\nu + 198\*\eta ^2\big)
+10\*\BS_{1,2}
+\BS_{2,1}\*\big(23-18\*\eta\big)
+\BS_{3}\*\big(36\*\eta-36\big)
+(N+1)\*N\*\Big(-\frct{2}{15}\*\BS_{2}\
+6\*\BS_{3} -3\*\BS_{2,1}\Big)
\Big]
\nonumber\\
&&{}
+\frct{64}{3}\*\frct{d_{AA}^{(4)}}{n_a}\*
  \Big[
 \frct{331}{15}
+ \frct{429}{2}\*\eta
- \frct{3824}{45}\*\nu
+ 180\*\eta^2
+ \BS_{1} \*\Big(-5 + 1352\*\eta - \frct{1328}{3}\*\nu + 722\*\eta^2 + 144\*\eta^3 + 192\*\nu^2\Big)
+ 2\*\BS_{2}\*\Big(\frct{1126}{15}+\frct{1127}{5}\*\eta
\nonumber\\
&&{}
- 19\*\nu + 48\*\nu^2 - 12\*\eta^2 \Big)
- 8\*\BS_{1, 1}\*(63\*\eta + 54\*\eta^2 - 25\*\nu +12\*\nu^2)
+ 9\*\BS_{3}\*(13 - 28\*\eta + 12\*\nu)
- \BS_{1, 2}\*(55 - 48\*\eta + 48\*\nu)
- 2\*\BS_{2, 1}\*(43
\nonumber\\
&&{}
- 66\*\eta + 30\*\nu)
+ 48\*\BS_{4}
- 54\*\BS_{1, 3}
- 18\*\BS_{2, 2}
- 24\*\BS_{3, 1}
+ 24\*\BS_{1, 1, 2}
+ 30\*\BS_{1, 2, 1}
- 6\*\BS_{2, 1, 1}
+ (N+1)\*N\*\Big(\frct{1}{15}\*\BS_{2}\
-6\*\BS_{3} +3\*\BS_{2,1}\Big)
\Big]


\,,\qquad
\label{eq:quartic}
\end{eqnarray}
\end{widetext}
where $\eta$ and $\nu$ are defined in Eq.~\eqref{eq:etanu} and we have omitted the arguments $N$ of the binomial harmonic sums.
The coefficients of the quadratic color factors require the full procedure described above.
To enable cancellations, we rewrite each binomial harmonic sum in the expression thus obtained in terms of nested harmonic sums.
We have
\begin{widetext}
\begin{eqnarray}
&&\gamma_{gg,\zeta_3}^{(3)}\Big|_{\mathrm{quad}}=





\frct{32}{3}\*C_F^3\*n_f\*
\Big[
24\*\D1
-24\*\D2
+231\*\D1^2
+294\*\D1^3
+144\*\D1^4
-\frct{897}{4}\*\eta
-16\*\nu
-177\*\eta^2
-138\*\eta^3
-114\*\eta^4
+6\*(
6\*\D1
\nonumber\\
&&{}
-6\*\D2
-2\*\D1^2
+8\*\D1^3
-71\*\eta
+20\*\nu
-78\*\eta^2
-56\*\eta^3
)\*\HS_{1}
-4\*(15\*\eta
-4\*\nu
+18\*\eta^2
)\*\HS_{2}
+4\*(33\*\eta
-16\*\nu
+90\*\eta^2
)\*\HS_{-2}
+12\*(13\*\eta
\nonumber\\
&&{}-4\*\nu
+14\*\eta^2
)\*\HS_{1,1}
-\frct{2}{3}\*\delta_{N,2}
\Big]

+\frct{32}{9}\*C_A\*C_F^2\*n_f\*\Big\{
-\frct{11}{2}
+\frct{1161}{2}\*(\D1-\D2)
-2817\*\D1^2
-1476\*\D1^3
-1836\*\D1^4
+148\*\D2^2
\nonumber\\
&&{}
+120\*\D2^3
+176\*\nu^2
+\frct{6879}{2}\*\eta
-\frct{1597}{3}\*\nu
+\frct{3771}{2}\*\eta^2
+2376\*\eta^3
+1224\*\eta^4
+3\*\Big[\frct{37}{2}
-30\*(\D1-\D2)
+123\*\D1^2
+12\*\D2^2
+132\*\D1^3
\nonumber\\
&&{}
+1065\*\eta
-366\*\nu
+1656\*\eta^2
-8\*\nu^2
+930\*\eta^3
\Big]\*\HS_{1}
+3\*(231\*\eta
-68\*\nu
+330\*\eta^2
)\*\HS_{2}
-24\*(24\*\eta
-13\*\nu
+81\*\eta^2
)\*\HS_{-2}
-48\*(30\*\eta
-10\*\nu
\nonumber\\
&&{}
+39\*\eta^2
)\*\HS_{1,1}
+\frct{7}{3}\*\delta_{N,2}
\Big\}

+\frct{32}{3}\*C_A^2\*C_F\*n_f\*\Big\{\frct{2}{3}
-\frct{471}{2}\*(\D1-\D2)
+\frct{2019}{2}\*\D1^2
-\frct{514}{3}\*\D2^2
+201\*\D1^3
-136\*\D2^3
+684\*\D1^4
\nonumber\\
&&{}
-\frct{2530}{3}\*\eta
+\frct{1303}{18}\*\nu
-\frct{1325}{3}\*\eta^2
-\frct{256}{3}\*\nu^2
-892\*\eta^3
-40\*\nu^3
-420\*\eta^4
+\Big[\frct{241}{6}
-6\*(\D1-\D2)
-135\*\D1^2
-84\*\D2^2
-324\*\D1^3
\nonumber\\
&&{}
-552\*\eta
+\frct{773}{3}\*\nu
-1470\*\eta^2
+112\*\nu^2
-726\*\eta^3
\Big]\*\HS_{1}
-(
309\*\eta
-124\*\nu
+426\*\eta^2
)\*\HS_{2}
+2\*(3
-3\*\eta
-8\*\nu
+210\*\eta^2
)\*\HS_{-2}
+8\*(60\*\eta
\nonumber\\
&&{}
-23\*\nu
+90\*\eta^2
)\*\HS_{1,1}
+8\*\HS_{3}
+8\*\HS_{-3}
-24\*\HS_{1,2}
-24\*\HS_{1,-2}
-24\*\HS_{2,1}
+8\*\HS_{-2,1}
-\frct{2}{9}\*\delta_{N,2}
\Big\}

+\frct{32}{3}\*C_A^3\*n_f\*\Big[
\frct{2197}{360}
-294\*\D1^2
+122\*\D2^2
\nonumber\\
&&{}
-24\*\D1^3
+96\*\D2^3
-216\*\D1^4
+\frct{683}{12}\*\eta
+\frct{6734}{135}\*\nu
+\frct{2431}{12}\*\eta^2
+\frct{194}{3}\*\nu^2
+377\*\eta^3
+20\*\nu^3
+144\*\eta^4
+\frct{1}{18}\*(
-1173
+432\*\D1^2
\nonumber\\
&&{}
+1296\*\D2^2
+2592\*\D1^3
+4989\*\eta
-1696\*\nu
+8913\*\eta^2
-1440\*\nu^2
+4194\*\eta^3
)\*\HS_{1}
+\frct{5}{3}\*(99\*\eta
-40\*\nu
+54\*\eta^2
)\*\HS_{2}
-\frct{2}{3}\*(
315\*\eta
-92\*\nu
\nonumber\\
&&{}
+162\*\eta^2
)\*\HS_{1,1}
+\frct{1}{6}\*\Big(
\frct{128}{15}
+\frct{2436}{5}\*\eta
-116\*\nu
-1056\*\eta^2
\Big)\*\HS_{-2}
-\frct{1}{6}\*(
59
-18\*\eta)\*\HS_{3}
+\frct{1}{3}\*(31
-18\*\eta)\*\HS_{-3}
+24\*\HS_{1,2}
+\frct{46}{3}\*\HS_{1,-2}
\nonumber\\
&&{}
+24\*\HS_{2,1}
+12\*(
-3
+\eta)\*\HS_{-2,1}
+\frct{1}{2}\*N\*(N+1)\*\Big(
\frct{4}{45}\*\HS_{-2}
+\HS_{3}
-2\*\HS_{-3}
+4\*\HS_{-2,1}\Big)\Big]

+\frct{64}{3}\*C_A^4\*\Big[
-\frct{2621}{720}
-\frct{36013}{144}\*\eta
+\frct{101027}{1080}\*\nu
\nonumber\\
&&{}
-\frct{1133}{6}\*\eta^2
-\frct{207}{4}\*\nu^2
-97\*\eta^3
+19\*\nu^3
-12\*\eta^4
-24\*\nu^4
+\frct{1}{36}\*(437
-21021\*\eta
+5915\*\nu
-12174\*\eta^2
-2664\*\nu^2
-4320\*\eta^3
\nonumber\\
&&{}
+1728\*\nu^3)\*\HS_{1}
+\frct{1}{3}\*(5
-51\*\eta
+33\*\nu)\*\HS_{3}
+6\*\HS_{4}
+\frct{4}{3}\*(45\*\eta
+45\*\eta^2
-25\*\nu
+12\*\nu^2)\*(\HS_{2}-2\*\HS_{1,1})
-\frct{1}{18}\*\Big(
\frct{452}{5}
-\frct{1023}{5}\*\eta
+435\*\nu
\nonumber\\
&&{}
-612\*\eta^2
+72\*\nu^2
\Big)\*\HS_{-2}
-\frct{1}{3}\*(
131
-141\*\eta
+105\*\nu)\*\HS_{-3}
-18\*\HS_{-4}
-11\*\HS_{1,3}
+\frct{1}{6}\*(143
-96\*\eta
+96\*\nu)\*\HS_{1,-2}
+(127
-156\*\eta
\nonumber\\
&&{}
+108\*\nu)\*\frct{1}{2}\*\HS_{-2,1}
+35\*\HS_{1,-3}
+8\*\HS_{2,-2}
-5\*\HS_{3,1}
+4\*\HS_{-2,-2}
+16\*\HS_{-2,2}
+43\*\HS_{-3,1}
-16\*\HS_{1,1,-2}
-54\*\HS_{1,-2,1}
-32\*\HS_{-2,1,1}
\nonumber\\
&&{}
-N\*(N+1)\*\Big(
\frct{1}{90}\*\HS_{-2}
+\frct{1}{4}\*\HS_{3}
-\frct{1}{2}\*\HS_{-3}
+\HS_{-2,1}
\Big)
\Big]


\,,
\label{eq:quadratic}
\end{eqnarray}
\end{widetext}
where $D_k$ is defined in Eq.~\eqref{eq:etanu} and we have omitted the arguments $N$ of the nested harmonic sums.
The $\zeta_3\delta_{N,2}$ terms with $\{C_F^2,C_AC_F,C_A^2\}C_Fn_f$ in Eq.~\eqref{eq:quadratic} are familiar from the $C_AC_Fn_f^2$ contribution to $\gamma_{gg,\zeta_3}^{(3)}$ \cite{Falcioni:2024qpd} and similar contributions to $\gamma_{\mathrm{ps}}^{(3)}$ with $C_AC_Fn_f^2$ \cite{Gehrmann:2023cqm} and $\gamma_{gq}^{(3)}$ with $C_AC_Fn_f^2$ \cite{Falcioni:2023tzp}.
Such terms do not yet appear in $\gamma_{ij}^{(2)}$ \cite{Larin:1993vu,Larin:1996wd,Moch:2004pa,Vogt:2004mw,Ablinger:2014nga,Blumlein:2021enk,Blumlein:2022gpp,Gehrmann:2023ksf}, but were already observed at two loops in the coefficient functions of inclusive deep-inelastic scattering (DIS) by photon exchange \cite{vanNeerven:1991nn,Moch:1999eb}.
They emerge from $D_{-2}S_{-2}(N-2)$ for $N=2$ and allow us to predict that $\gamma_{gg,\mathrm{rat}}^{(3)}$ will contain the term $(64/81)(6C_F^2-7C_AC_F+2C_A^2)C_Fn_f\theta(N-4)D_{-2}S_{-2}(N-2)$.
The $\zeta_3N(N+1)\BS_2$ terms with $n_fd_{RA}^{(4)}$ and $d_{AA}^{(4)}$ in Eq.~\eqref{eq:quartic} and with $n_fC_A^3$ and $C_A^4$ in Eq.~\eqref{eq:quadratic} (notice that $S_{-2}=-\BS_2/2$) conspire with similar terms in $\gamma_{gg,\zeta_5}^{(3)}$ \cite{Moch:2018wjh,Falcioni:2024qpd} to produce the first two terms of the well-known factor
\begin{equation}








f(N) = 5 \zeta_5 - 2 \zeta_3 \BS_{2}(N) - \BS_{2, 1, 2}(N) + \BS_{3, 1, 1}(N)\,,
\label{eq:fn}
\end{equation}
which was first encountered in three-loop coefficient functions of inclusive DIS \cite{Vermaseren:2005qc,Moch:2008fj} and later recovered in wrapping corrections to twist-two operators at four loops in $\mathcal{N}=4$ SYM \cite{Bajnok:2008qj}.
The analogous observation was made in Ref.~\cite{Falcioni:2024qpd} for the $\{C_A^2,d_{FF}^{(4)}/n_a\}n_f^2N(N+1)\{\zeta_5,\zeta_3\BS_{2}\}$ terms of $\gamma_{gg}^{(3)}$, where $d_{RR}^{(4)}/n_a=d_R^{abcd}d_R^{abcd}/n_a=(n_c^4-6n_c^2+18)(n_c^2-1)/(96n_c^3)$.
We thus conjecture that $\gamma_{gg,\mathrm{rat}}^{(3)}$ will contain the term $c_1N(N+1)(- \BS_{2, 1, 2} + \BS_{3, 1, 1})$, where $c_1=-(8/45)[C_A^2(C_A-n_f)^2/3+4(d_{AA}^{(4)}-4d_{RA}^{(4)}n_f+4d_{RR}^{(4)}n_f^2)/n_a]$.
A similar conjecture for $\gamma_{\mathrm{ns},\mathrm{rat}}^{(3)\mathrm{s}}$ \cite{Kniehl:2026axe} has recently been confirmed in Ref.~\cite{Gehrmann:2026qbl}.
In this way, these $N(N+1)$ power terms will be tamed to comply with the leading large-$N$ behavior of $\gamma_{gg}^{(3)}$ being $\ln N$ times the cusp anomalous dimension \cite{Korchemsky:1988si} because $f\propto\eta$ in this limit.
By the same token, the power divergence of the previously unobserved term $c_2N(N+1)(2\BS_3-\BS_{2,1})$, with $c_2=8[C_A^3(C_A-n_f)/3-8(d_{AA}^{(4)}-2d_{RA}^{(4)}n_f)/n_a]$, in $\gamma_{gg,\zeta_3}^{(3)}$ must be canceled by an appropriate counterpart in $\gamma_{gg,\mathrm{rat}}^{(3)}$, where we have used $2\BS_3-\BS_{2,1}=2(2S_{-3}-S_3-4S_{-2,1})$.
We note that $\mathcal{N}=1$ SYM is devoid of such terms because $c_1=c_2=0$ for $C_A=n_f$ and $d_{AA}^{(4)}=2d_{RA}^{(4)}n_f=4d_{RR}^{(4)}n_f^2$.
These color factor identifications are also encoded in Dokshitzer's relation among the entries in Eq.~\eqref{eq:matrix} \cite{Dokshitzer:1977sg,Bukhvostov:1985vj,Bukhvostov:1985rn,Belitsky:1998gu}, valid for structures appearing for the first time at a given order of perturbation theory, which we have used to transfer available information from $\gamma_{qq}^{(3)}$, $\gamma_{qg}^{(3)}$, and $\gamma_{gq}^{(3)}$ \cite{Falcioni:2023luc,Falcioni:2023vqq,Falcioni:2024xyt,Falcioni:2024qpd,Kniehl:2025jfs,Kniehl:2025ttz,Kniehl:2026axe,Gehrmann:2026qbl} to $\gamma_{gg}^{(3)}$.

In combination with Eq.~(8) in Ref.~\cite{Falcioni:2024qpd} and the $\zeta_3$ part of Eq.~(3.14) in Ref.~\cite{Davies:2016jie}, Eqs.~\eqref{eq:quartic} and \eqref{eq:quadratic} exhaust $\gamma_{gg,\zeta_3}^{(3)}$.
For the reader's convenience, we list the complete transcendental part of $\gamma_{gg}^{(3)}$, expressed in terms of nested harmonic sums, in the Supplemental Material~\ref{supplement:2}.
Apart from the $n_f^{1,0}$ terms of $\gamma_{qg,\zeta_3}^{(3)}$ and the $n_f^{0}$ term of $\gamma_{gq,\zeta_3}^{(3)}$, all transcendental contributions to $\gamma_{ij}^{(3)}$ are now available.

Having completed our knowledge of $\gamma_{gg,\zeta_3}^{(3)}$ and conjectured specific terms of $\gamma_{gg,\mathrm{rat}}^{(3)}$, we now present the $C_F^2n_f^2$ term of the latter for all $N$.
For this particular QED-like color factor, we meanwhile avail of $\gamma_{gg}^{(3)}$ through $N=48$, considerably beyond what is published \cite{Moch:2021qrk,Moch:2023tdj,Falcioni:2024qpd,Falcioni:2025hfz}.
The values for $N=2,\ldots,44$ turn out to suffice for analytic reconstruction via the LLL algorithm \cite{Lenstra82factoringpolynomials,fplll}, leaving those for $N=46,48$ for cross checking.
We find
\begin{widetext}
\begin{eqnarray}
&&\gamma_{gg,\mathrm{rat}}^{(3)}=





\frac{32}{27}\*C_F^2\*n_f^2\*\Big\{
-\frct{169}{16}
-\frct{1234}{3}\*\D1
+\frct{1234}{3}\*\D2
+\frct{63069}{32}\*\eta
+\frct{30839}{36}\*\nu
-\frct{30813}{8}\*\D1^2
-314\*\D2^2
-\frct{267373}{48}\*\eta^2
-\frct{2330}{27}\*\nu^2
\nonumber\\&&{}
-276\*\D2^3
-\frct{15611}{3}\*\D1^3
-\frct{55291}{12}\*\eta^3
-1566\*\D1^5
-111\*\D1^6
-540\*\D1^7
-128\*\D2^5
-3162\*\eta^5
-1476\*\eta^6
-270\*\eta^7
-\frct{640}{3}\*\D2^4
\nonumber\\&&{}
-\frct{21109}{4}\*\eta^4
-\frct{29527}{12}\*\D1^4
+
4\*(45\*\eta
+90\*\eta^2
+9\*\eta^3
-15\*\nu
-8\*\nu^2)\*\HS_{-3}
+\Big(
198\*\D2
-198\*\D1
+468\*\D1^2
+72\*\D1^3
+288\*\D1^4
\nonumber\\&&{}
-32\*\D2^3
+963\*\eta^2
+684\*\eta^3
+36\*\eta^4
-\frct{32}{3}\*\D2^2
+\frct{63}{2}\*\eta
-\frct{128}{3}\*\nu^2
-\frct{859}{9}\*\nu\Big)\*\HS_{-2}
+\Big[
\frct{785}{12}\*(\D1-\D2)
+\frct{161}{12}\*\D1^2
-40\*\D1^4
+12\*\D1^5
\nonumber\\&&{}
+388\*\eta^4
+57\*\eta^5
+\frct{1099}{6}\*\D1^3
+\frct{1213}{12}\*\eta^2
-\frct{1421}{8}\*\eta
+\frct{2257}{36}\*\nu
+\frct{9353}{12}\*\eta^3\Big]\*\HS_{1}
+\Big(
109\*\eta
-323\*\D1^2
+144\*\D1^3
-126\*\D1^4
+238\*\eta^3
\nonumber\\&&{}
+66\*\eta^4
+\frct{529}{6}\*\eta^2
-\frct{239}{18}\*\nu
+\frct{583}{2}\*\D1
-\frct{583}{2}\*\D2
\Big)\*\HS_{2}
+\Big(48\*\D1
-30\*\D1^2
+36\*\D1^3
-48\*\D2
-276\*\eta
-386\*\eta^2
-96\*\eta^3
-8\*\nu^2
\nonumber\\&&{}
+\frct{244}{3}\*\nu\Big)\*\HS_{3}
+2\*(27\*\eta
+12\*\eta^2
-14\*\nu)\*\HS_{4}
-6\*(27\*\eta
+48\*\eta^2
+24\*\eta^3
-8\*\nu)\*\HS_{-2,1}
-6\*(129\*\eta
+168\*\eta^2
+36\*\eta^3
-44\*\nu)\*\HS_{1,-2}
\nonumber\\&&{}
+\Big(33\*\D1^3
+105\*\D1^4
-231\*\eta^3
-87\*\eta^4
-\frct{45}{2}\*\D1
+\frct{45}{2}\*\D2
+\frct{471}{4}\*\D1^2
+\frct{579}{2}\*\eta
-\frct{2695}{18}\*\nu
+\frct{3943}{12}\*\eta^2\Big)\*\HS_{1,1}
+\Big(
-120\*\D1
+69\*\D1^2
\nonumber\\&&{}
-102\*\D1^3
+120\*\D2
+40\*\eta^2
-42\*\eta^3
-\frct{218}{3}\*\nu
+\frct{363}{2}\*\eta\Big)\*\HS_{1,2}
+2\*(27\*\eta
+42\*\eta^2
-4\*\nu)\*\HS_{1,3}
+\Big(60\*\D1
-30\*\D1^2
+60\*\D1^3
-60\*\D2
\nonumber\\&&{}
-216\*\eta^2
-3\*\eta^3
+138\*\nu
-\frct{1377}{4}\*\eta\Big)\*\HS_{1,1,1}
+\Big(
-96\*\D1
+57\*\D1^2
-78\*\D1^3
+96\*\D2
-46\*\eta^2
-42\*\eta^3
-\frct{46}{3}\*\nu
+\frct{81}{2}\*\eta\Big)\*\HS_{2,1}
\nonumber\\&&{}
-2\*(27\*\eta
+6\*\eta^2
-16\*\nu)\*\HS_{3,1}
+12\*(3\*\eta
+3\*\eta^2
-\nu)\*(\HS_{-4}
-2\*\HS_{-3,1}
-4\*\HS_{-2,-2}
-6\*\HS_{1,-3}
-4\*\HS_{2,-2}
+4\*\HS_{1,-2,1}
+8\*\HS_{1,1,-2})
\nonumber\\&&{}
+(
-9\*\eta
-6\*\eta^2
+4\*\nu)\*(11\*\HS_{2,2}
+3\*\HS_{1,1,2}
-\HS_{1,2,1}
-7\*\HS_{2,1,1}
-11\*\HS_{1,1,1,1})\Big\}


\,.
\label{eq:rat}
\end{eqnarray}
\end{widetext}
Equations~\eqref{eq:quartic}, \eqref{eq:quadratic}, and \eqref{eq:rat} represent our key results and mark an important step towards establishing exact $N$ dependencies of the $\gamma_{ij}^{(3)}$ set governing DGLAP evolution of PDFs at N${}^3$LO.

We briefly review our workflow.
To be specific, we consider $\gamma_{gg,\zeta_3}^{(3)}$.
We have $\gamma_{\mathrm{ns},\zeta_3}^{(3)+}$ \cite{Kniehl:2025jfs,Kniehl:2026axe,Gehrmann:2026qbl} and $\gamma_{\mathrm{ps},\zeta_3}^{(3)}$ \cite{Falcioni:2025hfz} for all $N$, and $\gamma_{gg,\zeta_3}^{(3)}$ \cite{Moch:2021qrk,Moch:2023tdj,Falcioni:2024qpd,Falcioni:2025hfz} and hence also $\tr\hat{\gamma}_{\zeta_3}^{(3)}=\gamma_{\mathrm{ns},\zeta_3}^{(3)+}+\gamma_{\mathrm{ps},\zeta_3}^{(3)}+\gamma_{gg,\zeta_3}^{(3)}$ for a set of low-$N$ values.
We evaluate the non-reciprocity-respecting component $\NP_{\zeta_3}^{(3)}$ of $\tr\hat{\gamma}_{\zeta_3}^{(3)}$ for all $N$ from sub-N${}^3$LO results via Eq.~(14) of Ref.~\cite{Kniehl:2026axe}.
We write $\mathcal{P}_{\zeta_3}^{(3)}=\tr\hat{\gamma}_{\zeta_3}^{(3)}+\NP_{\zeta_3}^{(3)}$ as a linear combination, with rational coefficients, of all binomial harmonic sums, bare or dressed by powers of $\eta$ or $\nu$, of overall weight $w\le4$.
We reduce the number of coefficients in this ansatz by imposing constraints from $\mathcal{N}=1,4$ SYM via Dokshitzer's relation \cite{Dokshitzer:1977sg,Bukhvostov:1985vj,Bukhvostov:1985rn,Belitsky:1998gu} and maximal-transcendentality principle \cite{Kotikov:2002ab,Kotikov:2006ts}, respectively.
Matching the low-$N$ values of $\mathcal{P}_{\zeta_3}^{(3)}$, we fix the coefficients of $\eta^{4,3,2}$ and $\nu^4$ using the Chinese remainder theorem \cite{china} and analytically reconstruct the remaining ones with \texttt{fplll} \cite{fplll}.
We end up with $\gamma_{gg,\zeta_3}^{(3)}=\mathcal{P}_{\zeta_3}^{(3)}-\NP_{\zeta_3}^{(3)}-\gamma_{\mathrm{ns},\zeta_3}^{(3)+}-\gamma_{\mathrm{ps},\zeta_3}^{(3)}$ for all $N$.
The evaluation of the $C_F^2n_f^2$ term of $\gamma_{gg,\mathrm{rat}}^{(3)}$ proceeds similarly.

We conclude with an outlook.
The present supply of analytic low-$N$ $\gamma_{ij}^{(3)}$ values, obtained with the FORCER program~\cite{Ruijl:2017cxj} written in FORM~\cite{Vermaseren:2000nd,Kuipers:2012rf,Ruijl:2017dtg}, needs further extension to enable analytic reconstruction of the as-yet missing all-$N$ expressions.
While this appears feasible for contributions involving $\zeta_3$ and/or $n_f$, non-fermionic rational contributions with $w=7$ may prove more difficult, as they require considerably more input.
Further progress may come from the alternative approach recently developed and applied to the QED-like $C_F^3n_f$ contributions to $\gamma_{\mathrm{ns}}^{(3)\pm}$ in Ref.~\cite{Gehrmann:2023iah} and to $\gamma_{\mathrm{ns},\mathrm{rat}}^{(3)\pm}$ and
$\gamma_{\mathrm{ns},\mathrm{rat}}^{(3)\mathrm{s}}$ in Ref.~\cite{Gehrmann:2026qbl}.
This relies on the computation of off-shell operator matrix elements and employs integration-by-parts reductions and differential equations with respect to a tracing parameter.
It is reassuring to notice that our recent LLL-based \cite{Lenstra82factoringpolynomials,fplll} analytic reconstructions of transcendental \cite{Kniehl:2025jfs,Kniehl:2026axe} and fermionic \cite{Kniehl:2025ttz} contributions to $\gamma_{\mathrm{ns}}^{(3)}$ were confirmed by Ref.~\cite{Gehrmann:2026qbl}.
The study of scaling violations in PDFs remains a challenging and intriguing topic of considerable phenomenological importance and broad interest, as the progress presented here exploits symmetries and structural relations in quantum field theory beyond QCD, also linking developments across different areas of theoretical physics.

\vspace{2mm}
\begin{acknowledgments}
\emph{Acknowledgments:}
B.A.K. and V.N.V. were supported in part by Deutsche Forschungsgemeinschaft (DFG) through Grants KN~365/13-2 and KN~365/16-1, S.O.M. by European Research Council (ERC) through Advanced Grant 101095857 {\it Conformal-EIC}, and A.V. by UK Science and Technology Facilities Council (STFC) through Consolidated Grant ST/X000699/1 and by Alexander von Humboldt-Stiftung (AvH) through a Research Award.
\end{acknowledgments}










\begingroup\begin{thebibliography}{99}


\bibitem{Bjorken:1969ja}
J.~D.~Bjorken and E.~A.~Paschos,
Inelastic Electron-Proton and $\gamma$-Proton Scattering and the Structure of the Nucleon,
Phys. Rev. \textbf{185}, 1975--1982 (1969)




\bibitem{Fritzsch:1973pi}
H.~Fritzsch, M.~Gell-Mann, and H.~Leutwyler,
Advantages of the color octet gluon picture,
Phys. Lett. B \textbf{47}, 365--368 (1973)




\bibitem{Gribov:1972ri}
  V.~N.~Gribov and L.~N.~Lipatov,
  Deep inelastic $ep$ scattering in perturbation theory,
  Yad.\ Fiz.\  {\bf 15}, 781--807 (1972)
  [Sov.\ J.\ Nucl.\ Phys.\  {\bf 15}, 438--450 (1972)].




\bibitem{Gribov:1972rt}
  V.~N.~Gribov and L.~N.~Lipatov,
  $e^+e^-$-pair annihilation and deep inelastic $ep$ scattering in perturbation theory,
  Yad.\ Fiz.\ {\bf 15}, 1218--1237 (1972)
  [Sov.\ J.\ Nucl.\ Phys.\ {\bf 15}, 675--684 (1972)].




\bibitem{Altarelli:1977zs}
  G.~Altarelli and G.~Parisi,
  Asymptotic freedom in parton language,
  Nucl.\ Phys.\ B {\bf 126}, 298--318 (1977).





\bibitem{Dokshitzer:1977sg}
  Yu.~L.~Dokshitzer,
  Calculation of structure functions of deep-inelastic scattering and $e^+ e^-$ annihilation by perturbation theory in quantum chromodynamics,
  Zh.\ Eksp.\ Teor.\ Fiz.\  {\bf 73}, 1216--1240 (1977)
  [Sov.\ Phys.\ JETP {\bf 46}, 641--653 (1977)].




\bibitem{Collins:1989gx}
J.~C.~Collins, D.~E.~Soper, and G.~Sterman,
Factorization of Hard Processes in QCD,
Adv. Ser. Direct. High Energy Phys. \textbf{5}, 1--91 (1989)

[arXiv:hep-ph/0409313 [hep-ph]].



\bibitem{Alekhin:2025qdj}
S.~Alekhin, M.~V.~Garzelli, S.~Moch, and O.~Zenaiev,
Determination of the strong coupling from high-energy data,
Eur. Phys. J. C \textbf{86} (2026) 161

[arXiv:2510.21435 [hep-ph]].



\bibitem{Alekhin:2010sv}
S.~Alekhin and S.~Moch,
Heavy-quark deep-inelastic scattering with a running mass,
Phys. Lett. B \textbf{699} (2011) 345--353

[arXiv:1011.5790 [hep-ph]].



\bibitem{Alekhin:2012un}
S.~Alekhin, K.~Daum, K.~Lipka and S.~Moch,
Determination of the charm-quark mass in the $\overline{\mathrm{MS}}$ scheme using charm production data from deep inelastic scattering at HERA,
Phys. Lett. B \textbf{718} (2012) 550--557

[arXiv:1209.0436 [hep-ph]].



\bibitem{Gross:1973ju}
D.~J.~Gross and F.~Wilczek,
Asymptotically Free Gauge Theories. I,
Phys. Rev. D \textbf{8}, 3633--3652 (1973)




\bibitem{Gross:1974cs}
D.~J.~Gross and F.~Wilczek,
Asymptotically free gauge theories. II,
Phys. Rev. D \textbf{9}, 980--993 (1974)




\bibitem{Floratos:1977au}
  E.~G.~Floratos, D.~A.~Ross, and C.~T.~Sachrajda,
  Higher-order effects in asymptotically free gauge theories: The anomalous dimensions of Wilson operators,
  Nucl.\ Phys.\ B {\bf 129}, 66--88 (1977); {\bf 139}, 545--546(E) (1978)







\bibitem{GonzalezArroyo:1979df}
  A.~Gonz\'alez-Arroyo, C.~L\'opez, and F.~J.~Yndur\'ain,
  Second-order contributions to the structure functions in deep inelastic scattering (I). Theoretical calculations,
  Nucl.\ Phys.\ B {\bf 153}, 161--186 (1979)





\bibitem{Gonzalez-Arroyo:1979kjx}
A.~Gonz\'alez-Arroyo, C.~L\'opez, and F.~J.~Yndur\'ain,
Testing QCD predictions beyond the leading order for non-singlet structure functions in neutrino scattering,
Nucl.\ Phys.\ B \textbf{174}, 474--490 (1980)




\bibitem{Curci:1980uw}
G.~Curci, W.~Furmanski, and R.~Petronzio,
Evolution of parton densities beyond leading order: The non-singlet case,
Nucl.\ Phys.\ B \textbf{175}, 27--92 (1980)




\bibitem{Furmanski:1980cm}
W.~Furmanski and R.~Petronzio,
Singlet parton densities beyond leading order,
Phys.\ Lett.\ B \textbf{97}, 437--442 (1980)




\bibitem{Larin:1993vu}
S.~A.~Larin, T.~van Ritbergen, and J.~A.~M.~Vermaseren,
The next-next-to-leading QCD approximation for non-singlet moments of deep inelastic structure functions,''
Nucl.\ Phys.\ B \textbf{427}, 41--52 (1994)




\bibitem{Larin:1996wd}
S.~A.~Larin, P.~Nogueira, T.~van Ritbergen, and J.~A.~M.~Vermaseren,
The 3-loop QCD calculation of the moments of deep inelastic structure functions,
Nucl.\ Phys.\ B \textbf{492}, 338--378 (1997)

[arXiv:hep-ph/9605317 [hep-ph]].



\bibitem{Moch:2004pa}
S.~Moch, J.~A.~M.~Vermaseren, and A.~Vogt,
The three-loop splitting functions in QCD: the non-singlet case,
Nucl.\ Phys.\ B {\bf 688}, 101--134 (2004)

[hep-ph/0403192].




\bibitem{Vogt:2004mw}
A.~Vogt, S.~Moch, and J.~A.~M.~Vermaseren,
The three-loop splitting functions in QCD: the singlet case,
Nucl.\ Phys.\ B {\bf 691}, 129--181 (2004)

[hep-ph/0404111].




\bibitem{Ablinger:2014nga}
J.~Ablinger, A.~Behring, J.~Bl{\"u}mlein, A.~De Freitas, A.~von Manteuffel, and C.~Schneider,
The 3-loop pure singlet heavy flavor contributions to the structure function $F_2(x,Q^2)$ and the anomalous dimension,
Nucl.\ Phys.\ B \textbf{890}, 48--151 (2014)

[arXiv:1409.1135 [hep-ph]].



\bibitem{Blumlein:2021enk}
J.~Bl{\"u}mlein, P.~Marquard, C.~Schneider, and K.~Sch{\"o}nwald,
The three-loop unpolarized and polarized non-singlet anomalous dimensions from off shell operator matrix elements,
Nucl.\ Phys.\ B \textbf{971}, 115542 (2021)

[arXiv:2107.06267 [hep-ph]].



\bibitem{Blumlein:2022gpp}
J.~Bl{\"u}mlein, P.~Marquard, C.~Schneider, and K.~Sch{\"o}nwald,
The massless three-loop Wilson coefficients for the deep-inelastic structure functions $F_{2}$, $F_{L}$, $xF_{3}$ and $g_{1}$,
JHEP \textbf{11}, 156 (2022)

[arXiv:2208.14325 [hep-ph]].



\bibitem{Gehrmann:2023ksf}
T.~Gehrmann, A.~von Manteuffel, and T.~Z.~Yang,
Renormalization of twist-two operators in covariant gauge to three loops in QCD,
JHEP \textbf{04}, 041 (2023)

[arXiv:2302.00022 [hep-ph]].



\bibitem{Gracey:1994nn}
J.~A.~Gracey,
Anomalous dimension of non-singlet Wilson operators at $O(1 / N_f)$ in deep inelastic scattering,
Phys. Lett. B \textbf{322}, 141--146 (1994)

[arXiv:hep-ph/9401214 [hep-ph]].



\bibitem{Velizhanin:2011es}
V.~N.~Velizhanin,
Four loop anomalous dimension of the second moment of the non-singlet twist-2 operator in QCD,
Nucl. Phys. B \textbf{860}, 288--294 (2012)

[arXiv:1112.3954 [hep-ph]].



\bibitem{Velizhanin:2014fua}
V.~N.~Velizhanin,
Four-loop anomalous dimension of the third and fourth moments of the nonsinglet twist-2 operator in QCD,
Int. J. Mod. Phys. A \textbf{35}, 2050199 (2020)

[arXiv:1411.1331 [hep-ph]].



\bibitem{Davies:2016jie}
J.~Davies, A.~Vogt, B.~Ruijl, T.~Ueda, and J.~A.~M.~Vermaseren,
Large-$n_f$ contributions to the four-loop splitting functions in QCD,
Nucl. Phys. B \textbf{915}, 335--362 (2017)

[arXiv:1610.07477 [hep-ph]].



\bibitem{Moch:2017uml}
S.~Moch, B.~Ruijl, T.~Ueda, J.~A.~M.~Vermaseren, and A.~Vogt,
Four-loop non-singlet splitting functions in the planar limit and beyond,
JHEP \textbf{10}, 041 (2017)

[arXiv:1707.08315 [hep-ph]].



\bibitem{Davies:2017hyl}
J.~Davies and A.~Vogt,
Absence of $\pi^2$ terms in physical anomalous dimensions in DIS: Verification and resulting predictions,
Phys. Lett. B \textbf{776}, 189--194 (2018)

[arXiv:1711.05267 [hep-ph]].



\bibitem{Moch:2018wjh}
S.~Moch, B.~Ruijl, T.~Ueda, J.~A.~M.~Vermaseren, and A.~Vogt,
On quartic colour factors in splitting functions and the gluon cusp anomalous dimension,
Phys. Lett. B \textbf{782}, 627--632 (2018)

[arXiv:1805.09638 [hep-ph]].



\bibitem{Das:2020adl}
G.~Das, S.~Moch, and A.~Vogt,
Approximate four-loop QCD corrections to the Higgs-boson production cross section,
Phys. Lett. B \textbf{807}, 135546 (2020)

[arXiv:2004.00563 [hep-ph]].



\bibitem{Moch:2021qrk}
S.~Moch, B.~Ruijl, T.~Ueda, J.~A.~M.~Vermaseren, and A.~Vogt,
Low moments of the four-loop splitting functions in QCD,
Phys. Lett. B \textbf{825}, 136853 (2022)

[arXiv:2111.15561 [hep-ph]].



\bibitem{Falcioni:2023luc}
G.~Falcioni, F.~Herzog, S.~Moch, and A.~Vogt,
Four-loop splitting functions in QCD -- The quark-quark case,
Phys. Lett. B \textbf{842}, 137944 (2023)

[arXiv:2302.07593 [hep-ph]].



\bibitem{Gehrmann:2023cqm}
T.~Gehrmann, A.~von Manteuffel, V.~Sotnikov, and T.-Z.~Yang,
Complete ${N}_f^2$ contributions to four-loop pure-singlet splitting functions,
JHEP \textbf{01}, 029 (2024)

[arXiv:2308.07958 [hep-ph]].



\bibitem{Falcioni:2023vqq}
G.~Falcioni, F.~Herzog, S.~Moch, and A.~Vogt,
Four-loop splitting functions in QCD -- The gluon-to-quark case,
Phys. Lett. B \textbf{846}, 138215 (2023)

[arXiv:2307.04158 [hep-ph]].



\bibitem{Falcioni:2023tzp}
G.~Falcioni, F.~Herzog, S.~Moch, J.~Vermaseren, and A.~Vogt,
The double fermionic contribution to the four-loop quark-to-gluon splitting function,
Phys. Lett. B \textbf{848}, 138351 (2024)

[arXiv:2310.01245 [hep-ph]].



\bibitem{Moch:2023tdj}
S.~Moch, B.~Ruijl, T.~Ueda, J.~Vermaseren, and A.~Vogt,
Additional moments and $x$-space approximations of four-loop splitting functions in QCD,
Phys. Lett. B \textbf{849}, 138468 (2024)

[arXiv:2310.05744 [hep-ph]].



\bibitem{Gehrmann:2023iah}
T.~Gehrmann, A.~von Manteuffel, V.~Sotnikov, and T.-Z.~Yang,
The $N_fC_F^3$ contribution to the non-singlet splitting function at four-loop order,
Phys. Lett. B \textbf{849}, 138427 (2024)

[arXiv:2310.12240 [hep-ph]].



\bibitem{Falcioni:2024xyt}
G.~Falcioni, F.~Herzog, S.~Moch, A.~Pelloni, and A.~Vogt,
Four-loop splitting functions in QCD -- The quark-to-gluon case,
Phys. Lett. B \textbf{856}, 138906 (2024)

[arXiv:2404.09701 [hep-ph]].



\bibitem{Falcioni:2024xav}
G.~Falcioni, F.~Herzog, S.~Moch, and S.~Van Thurenhout,
Constraints for twist-two alien operators in QCD,
JHEP \textbf{11}, 080 (2024)

[arXiv:2409.02870 [hep-ph]].



\bibitem{Falcioni:2024qpd}
G.~Falcioni, F.~Herzog, S.~Moch, A.~Pelloni, and A.~Vogt,
Four-loop splitting functions in QCD -- the gluon-gluon case --,
Phys. Lett. B \textbf{860}, 139194 (2025)

[arXiv:2410.08089 [hep-ph]].



\bibitem{Kniehl:2025jfs}
B.~A.~Kniehl and V.~N.~Velizhanin,
Four-Loop Anomalous Dimension of Flavor Nonsinglet Twist-Two Operator of General Lorentz Spin in QCD: $\zeta(3)$ Term,
Phys. Rev. Lett. \textbf{134}, 131901 (2025)

[arXiv:2503.20422 [hep-ph]].



\bibitem{Kniehl:2025ttz}
B.~A.~Kniehl, S.~Moch, V.~N.~Velizhanin, and A.~Vogt,
Flavor Nonsinglet Splitting Functions at Four Loops in QCD: Fermionic Contributions,
Phys. Rev. Lett. \textbf{135}, 071902 (2025)

[arXiv:2505.09381 [hep-ph]].



\bibitem{Falcioni:2025hfz}
G.~Falcioni, F.~Herzog, S.~Moch, A.~Pelloni, and A.~Vogt,
Additional results on the four-loop flavour-singlet splitting functions in QCD,
Phys.\ Lett.\ B \textbf{875}, 140278 (2026)

[arXiv:2512.10783 [hep-ph]].



\bibitem{Kniehl:2026axe}
B.~A.~Kniehl and V.~N.~Velizhanin,
Four-loop anomalous dimension of flavor non-singlet quark operator of twist two and Lorentz spin $N$ for general gauge group: Transcendental part,
Nucl.\ Phys.\ B \textbf{1028}, 117492 (2026)

[arXiv:2605.10284 [hep-ph]].



\bibitem{Gehrmann:2026qbl}
T.~Gehrmann, A.~von Manteuffel, V.~Sotnikov, and T.~Yang,
The four-loop non-singlet splitting functions in QCD,
[arXiv:2604.09534 [hep-ph]].



\bibitem{Jamin:2017mul}
M.~Jamin and R.~Miravitllas,
Absence of even-integer $\zeta$-function values in Euclidean physical quantities in QCD,
Phys. Lett. B \textbf{779}, 452--455 (2018)

[arXiv:1711.00787 [hep-ph]].



\bibitem{Baikov:2018wgs}
P.~A.~Baikov and K.~G.~Chetyrkin,
The structure of generic anomalous dimensions and no-$\pi$ theorem for massless propagators,
JHEP \textbf{06}, 141 (2018)

[arXiv:1804.10088 [hep-ph]].



\bibitem{Vermaseren:1998uu}
J.~A.~M.~Vermaseren,
Harmonic sums, Mellin transforms and integrals,
Int. J. Mod. Phys. A \textbf{14} (1999) 2037--2076

[arXiv:hep-ph/9806280 [hep-ph]].


\bibitem{Lenstra82factoringpolynomials}
A.~K.~Lenstra, H.~W.~Lenstra, and L.~Lov\'asz,
Factoring polynomials with rational coefficients,
Math.\ Ann.\ {\bf 261}, 515--534 (1982)



\bibitem{fplll}
M.~Albrech, D.~Cad\'{e}, X.~Pujol, and D.~Stehl\'{e} (The {FPLLL} development team),
\texttt{fplll}, a lattice reduction library, 2016,
available at {\tt https://github.com/fplll/fplll}.


\bibitem{Velizhanin:2010cm}
V.~N.~Velizhanin,
Six-loop anomalous dimension of twist-three operators in $\mathcal{N}=4$ SYM,
JHEP \textbf{11}, 129 (2010)

[arXiv:1003.4717 [hep-th]].



\bibitem{Velizhanin:2013vla}
V.~N.~Velizhanin,
Twist-2 at five loops: wrapping corrections without wrapping computations,
JHEP \textbf{06}, 108 (2014)

[arXiv:1311.6953 [hep-th]].



\bibitem{Marboe:2014sya}
C.~Marboe, V.~Velizhanin, and D.~Volin,
Six-loop anomalous dimension of twist-two operators in planar $\mathcal{N}=4$ SYM theory,
JHEP \textbf{07}, 084 (2015)

[arXiv:1412.4762 [hep-th]].



\bibitem{Marboe:2016igj}
C.~Marboe and V.~Velizhanin,
Twist-2 at seven loops in planar $\mathcal{N}=4$ SYM theory: full result and analytic properties,
JHEP \textbf{11}, 013 (2016)

[arXiv:1607.06047 [hep-th]].



\bibitem{Kniehl:2020rip}
B.~A.~Kniehl and V.~N.~Velizhanin,
Nonplanar cusp and transcendental anomalous dimension at four loops in $\mathcal{N}=4$ supersymmetric Yang-Mills theory,
Phys.\ Rev.\ Lett.\ \textbf{126}, 061603 (2021)

[arXiv:2010.13772 [hep-th]].



\bibitem{Kniehl:2021ysp}
B.~A.~Kniehl and V.~N.~Velizhanin,
Non-planar universal anomalous dimension of twist-two operators with general Lorentz spin at four loops in $\mathcal{N}=4$ SYM theory,
Nucl.\ Phys.\ B \textbf{968}, 115429 (2021)

[arXiv:2103.16420 [hep-th]].



\bibitem{Velizhanin:2021bdh}
V.~N.~Velizhanin,
NNNLLA BFKL pomeron eigenvalue in the planar $\mathcal{N}=4$ SYM theory,
[arXiv:2106.06527 [hep-th]].



\bibitem{Kniehl:2023bbk}
B.~A.~Kniehl and V.~N.~Velizhanin,
Anomalous dimensions of twist-two operators in extended $\mathcal{N}=2$ and $\mathcal{N}=4$ super Yang-Mills theories,
Nucl. Phys. B \textbf{1001}, 116511 (2024)

[arXiv:2312.05888 [hep-th]].



\bibitem{Kniehl:2024tvd}
B.~A.~Kniehl and V.~N.~Velizhanin,
Nonplanar four-loop anomalous dimensions of twist-two operators in $\mathcal{N}=4$ supersymmetric Yang-Mills theory: Higher moment, general result, and cusp anomalous dimension,
Phys. Rev. D \textbf{111}, L061902 (2025)

[arXiv:2409.09463 [hep-th]].



\bibitem{Velizhanin:2012nm}
V.~N.~Velizhanin,
Three loop anomalous dimension of the non-singlet transversity operator in QCD,
Nucl. Phys. B \textbf{864}, 113--140 (2012)

[arXiv:1203.1022 [hep-ph]].


\bibitem{china} Sun Tzu Suan Ching, 3rd to 5th century AD, see, {\it e.g.}, D.~E.~Knuth, The Art of Computer Programming, vol.~1: Fundamental Algorithms, 3rd ed., Adison--Wesley, Reading, Massachusetts, 1997.


\bibitem{Moch:2014sna}
S.~Moch, J.~A.~M.~Vermaseren, and A.~Vogt,
The three-loop splitting functions in QCD: The helicity-dependent case,
Nucl. Phys. B \textbf{889}, 351--400 (2014)

[arXiv:1409.5131 [hep-ph]].



\bibitem{Drell:1969jm}
S.~D.~Drell, D.~J.~Levy, and T.-M.~Yan,
Theory of Deep-Inelastic Lepton-Nucleon Scattering and Lepton-Pair Annihilation Processes.~I,
Phys. Rev. \textbf{187}, 2159--2171 (1969)



\bibitem{Blumlein:2000wh}
J.~Bl\"umlein, V.~Ravindran, and W.~L.~van Neerven,
On the Drell--Levy--Yan relation to $O(\alpha_s^2)$,
Nucl.\ Phys.\ B \textbf{586}, 349--381 (2000)

[arXiv:hep-ph/0004172 [hep-ph]].


\bibitem{Dokshitzer:2005bf}
Yu.~L.~Dokshitzer, G.~Marchesini, and G.~P.~Salam,
Revisiting parton evolution and the large-$x$ limit,
Phys. Lett. B \textbf{634}, 504--507 (2006)

[arXiv:hep-ph/0511302 [hep-ph]].



\bibitem{Dokshitzer:2006nm}
Yu.~L.~Dokshitzer and G.~Marchesini,
$\mathcal{N}=4$ SUSY Yang-Mills: Three loops made simple(r),
Phys. Lett. B \textbf{646}, 189--201 (2007)

[arXiv:hep-th/0612248 [hep-th]].



\bibitem{Basso:2006nk}
B.~Basso and G.~P.~Korchemsky,
Anomalous dimensions of high-spin operators beyond the leading order,
Nucl. Phys. B \textbf{775}, 1--30 (2007)

[arXiv:hep-th/0612247 [hep-th]].



\bibitem{Gross:1973id}
D.~J.~Gross and F.~Wilczek,
Ultraviolet Behavior of Non-Abelian Gauge Theories,
Phys. Rev. Lett. \textbf{30} (1973) 1343--1346




\bibitem{Politzer:1973fx}
H.~D.~Politzer,
Reliable Perturbative Results for Strong Interactions?,
Phys. Rev. Lett. \textbf{30} (1973) 1346--1349




\bibitem{Jones:1974mm}
D.~R.~T.~Jones,
Two-loop diagrams in Yang-Mills theory,
Nucl. Phys. B \textbf{75} (1974) 531--538




\bibitem{Caswell:1974gg}
W.~E.~Caswell,
Asymptotic Behavior of Non-Abelian Gauge Theories to Two-Loop Order,
Phys. Rev. Lett. \textbf{33} (1974) 244--246




\bibitem{Tarasov:1980au}
O.~V.~Tarasov, A.~A.~Vladimirov, and A.~Yu.~Zharkov,
The Gell-Mann-Low function of QCD in the three-loop approximation,
Phys. Lett. B \textbf{93} (1980) 429--432




\bibitem{Larin:1993tp}
S.~A.~Larin and J.~A.~M.~Vermaseren,
The three-loop QCD $\beta$-function and anomalous dimensions,
Phys. Lett. B \textbf{303} (1993) 334--336

[arXiv:hep-ph/9302208 [hep-ph]].



\bibitem{vanRitbergen:1997va}
T.~van Ritbergen, J.~A.~M.~Vermaseren, and S.~A.~Larin,
The four-loop $\beta$-function in quantum chromodynamics,
Phys. Lett. B \textbf{400}, 379--384 (1997)

[arXiv:hep-ph/9701390 [hep-ph]].



\bibitem{Chetyrkin:2004mf}
K.~G.~Chetyrkin,
Four-loop renormalization of QCD: Full set of renormalization constants and anomalous dimensions,
Nucl. Phys. B \textbf{710}, 499--510 (2005)

[arXiv:hep-ph/0405193 [hep-ph]].



\bibitem{Czakon:2004bu}
M.~Czakon,
The four-loop QCD $\beta$-function and anomalous dimensions,
Nucl. Phys. B \textbf{710}, 485--498 (2005)

[arXiv:hep-ph/0411261 [hep-ph]].



\bibitem{Luthe:2016ima}
T.~Luthe, A.~Maier, P.~Marquard, and Y.~Schr{\"o}der,
Towards the five-loop Beta function for a general gauge group,
JHEP \textbf{07}, 127 (2016)

[arXiv:1606.08662 [hep-ph]].



\bibitem{Baikov:2016tgj}
P.~A.~Baikov, K.~G.~Chetyrkin, and J.~H.~K{\"u}hn,
Five-Loop Running of the QCD Coupling Constant,
Phys. Rev. Lett. \textbf{118}, 082002 (2017)

[arXiv:1606.08659 [hep-ph]].



\bibitem{Herzog:2017ohr}
F.~Herzog, B.~Ruijl, T.~Ueda, J.~A.~M.~Vermaseren, and A.~Vogt,
The five-loop beta function of Yang-Mills theory with fermions,
JHEP \textbf{02}, 090 (2017)

[arXiv:1701.01404 [hep-ph]].



\bibitem{Luthe:2017ttg}
T.~Luthe, A.~Maier, P.~Marquard, and Y.~Schr{\"o}der,
The five-loop Beta function for a general gauge group and anomalous dimensions beyond Feynman gauge,
JHEP \textbf{10}, 166 (2017)

[arXiv:1709.07718 [hep-ph]].


\bibitem{Alday:2015eya}
L.~F.~Alday, A.~Bissi, and T.~{\L}ukowski,
Large spin systematics in CFT,
JHEP \textbf{11}, 101 (2015)

[arXiv:1502.07707 [hep-th]].


\bibitem{Chen:2020uvt}
H.~Chen, T.-Z.~Yang, H.-X.~Zhu, and Y.-J.~Zhu,
Analytic continuation and reciprocity relation for collinear splitting in QCD,
Chin.\ Phys.\ C \textbf{45}, 043101 (2021)

[arXiv:2006.10534 [hep-ph]].


\bibitem{Moch-Shumilov-to-appear}
S.-O.~Moch and L.~A.~Shumilov,
Analytic properties of anomalous dimensions of twist-two composite operators,
in preparation.


\bibitem{Kotikov:2002ab}
A.~V.~Kotikov and L.~N.~Lipatov,
DGLAP and BFKL equations in the $N=4$ supersymmetric gauge theory,
Nucl. Phys. B \textbf{661}, 19--61 (2003);
\textbf{685}, 405--407(E) (2004)


[arXiv:hep-ph/0208220 [hep-ph]].



\bibitem{Kotikov:2006ts}
A.~V.~Kotikov and L.~N.~Lipatov,
On the highest transcendentality in $\mathcal{N}=4$ SUSY,
Nucl. Phys. B \textbf{769}, 217-255 (2007)

[arXiv:hep-th/0611204 [hep-th]].



\bibitem{Beisert:2006ez}
N.~Beisert, B.~Eden, and M.~Staudacher,
Transcendentality and crossing,
J. Stat. Mech. \textbf{0701}, P01021 (2007)

[arXiv:hep-th/0610251 [hep-th]].



\bibitem{Bajnok:2008qj}
Z.~Bajnok, R.~A.~Janik, and T.~{\L}ukowski,
Four loop twist two, BFKL, wrapping and strings,
Nucl. Phys. B \textbf{816}, 376--398 (2009)

[arXiv:0811.4448 [hep-th]].



\bibitem{vanNeerven:1991nn}
W.~L.~van Neerven and E.~B.~Zijlstra,
Order $\alpha_s^2$ contributions to the deep inelastic Wilson coefficient,
Phys. Lett. B \textbf{272}, 127--133 (1991)




\bibitem{Moch:1999eb}
S.~Moch and J.~A.~M.~Vermaseren,
Deep-inelastic structure functions at two loops,
Nucl. Phys. B \textbf{573}, 853--907 (2000)

[arXiv:hep-ph/9912355 [hep-ph]].



\bibitem{Vermaseren:2005qc}
J.~A.~M.~Vermaseren, A.~Vogt, and S.~Moch,
The third-order QCD corrections to deep-inelastic scattering by photon exchange,
Nucl. Phys. B \textbf{724}, 3--182 (2005)

[arXiv:hep-ph/0504242 [hep-ph]].



\bibitem{Moch:2008fj}
S.~Moch, J.~A.~M.~Vermaseren, and A.~Vogt,
Third-order QCD corrections to the charged-current structure function $F_3$,
Nucl. Phys. B \textbf{813}, 220--258 (2009)

[arXiv:0812.4168 [hep-ph]].



\bibitem{Korchemsky:1988si}
G.~P.~Korchemsky,
Asymptotics of the Altarelli-Parisi-Lipatov Evolution Kernels of Parton Distributions,
Mod. Phys. Lett. A \textbf{4}, 1257--1276 (1989)




\bibitem{Bukhvostov:1985vj}
A.~P.~Bukhvostov, \'E.~A.~Kuraev, L.~N.~Lipatov, and G.~V.~Frolov,
Supersymmetry and anomalous dimensionalities of quasiparton operators in quantum chromodynamics,

Pis'ma Zh. Eksp. Teor. Fiz. \textbf{41}, 77--79  (1985)
[JETP Lett.\ \textbf{41}, 92--95 (1985)].



\bibitem{Bukhvostov:1985rn}
A.~P.~Bukhvostov, G.~V.~Frolov, L.~N.~Lipatov, and E.~A.~Kuraev,
Evolution equations for quasi-partonic operators,
Nucl.\ Phys.\ B \textbf{258}, 601--646 (1985)



\bibitem{Belitsky:1998gu}
A.~V.~Belitsky, D.~M\"uller, and A.~Sch\"afer,
Implications of $N=1$ supersymmetry for QCD conformal operators,
Phys. Lett. B \textbf{450}, 126--135 (1999)

[arXiv:hep-ph/9811484 [hep-ph]].



















\bibitem{Ruijl:2017cxj}
B.~Ruijl, T.~Ueda and J.~A.~M.~Vermaseren,
FORCER, a FORM program for the parametric reduction of four-loop massless propagator diagrams,
Comput. Phys. Commun. \textbf{253}, 107198 (2020)

[arXiv:1704.06650 [hep-ph]].



\bibitem{Vermaseren:2000nd}
J.~A.~M.~Vermaseren,
New features of FORM,
[arXiv:math-ph/0010025 [math-ph]].



\bibitem{Kuipers:2012rf}
J.~Kuipers, T.~Ueda, J.~A.~M.~Vermaseren, and J.~Vollinga,
FORM version 4.0,
Comput. Phys. Commun. \textbf{184}, 1453--1467 (2013)

[arXiv:1203.6543 [cs.SC]].



\bibitem{Ruijl:2017dtg}
B.~Ruijl, T.~Ueda, and J.~A.~M~Vermaseren,
FORM version 4.2,
[arXiv:1707.06453 [hep-ph]].


\end{thebibliography}\endgroup



\newpage
\cleardoublepage
\appendix



\renewcommand{\theequation}{A.\arabic{equation}}
\setcounter{equation}{0}
\renewcommand{\thefigure}{A.\arabic{figure}}
\setcounter{figure}{0}
\renewcommand{\thetable}{A.\arabic{table}}
\setcounter{table}{0}

\begin{widetext}
\setcounter{secnumdepth}{2}
\section{Supplemental Material}


\subsection{Typical application of \texttt{\textbf{fplll}}}
\label{supplement:1}

As a typical example for the analytic reconstruction using the program \texttt{fplll} \cite{fplll}, we consider the coefficient of $C_F^3\nf$ in $\gamma^{(3)}_{gg,\zeta_3}$.
We write down the full basis of binomial harmonic sums, bare and dressed by powers of $\eta$ and $\nu$, with overall weights $w\le4$,
\begin{eqnarray}
&&\{1, \eta, \nu, \eta^2, \nu^2, \BS_{1}, \eta \BS_{1},
 \nu \BS_{1}, \eta^3, \nu^3, \eta^2 \BS_{1}, \nu^2 \BS_{1}, \BS_{2},
 \eta \BS_{2}, \nu \BS_{2}, \BS_{1, 1}, \eta \BS_{1, 1},
 \nu \BS_{1, 1}, \eta^4, \nu^4, \eta^3 \BS_{1}, \nu^3 \BS_{1},
 \eta^2 \BS_{2}, \nu^2 \BS_{2}, \BS_{3}, \nonumber\\&&\ \ \
 \eta \BS_{3}, \nu \BS_{3}, \eta^2 \BS_{1, 1}, \nu^2 \BS_{1, 1}, \BS_{1, 2}, \eta \BS_{1, 2},
 \nu \BS_{1, 2}, \BS_{2, 1}, \eta \BS_{2, 1}, \nu \BS_{2, 1}\}\,.
 \label{BasisCF3nf}
\end{eqnarray}
It is useful to rescale the basis elements in Eq.~\eqref{BasisCF3nf} by a common factor $2^i/3^j$ to make their coefficients come out as small numbers.
Here, $32/3$ is found to yield optimal results.

We use as inputs the nine even Mellin moments with $N=4,6,\ldots,20$ \cite{Moch:2021qrk,Moch:2023tdj,Falcioni:2024qpd}, leaving the one with $N=22$ \cite{Falcioni:2025hfz} for cross checking.
We refrain from inputting the value for $N=2$ \cite{Moch:2021qrk}, bearing in mind that it might be contaminated by $\delta_{N,2}$, similarly to the $C_AC_Fn_f^2$ contribution to $\gamma_{gg,\zeta_3}^{(3)}$ \cite{Falcioni:2024qpd}.
Applying \texttt{fplll} with parameters ``\texttt{-f mpfr -p 100 -a bkz -b 10}'' to the matrix thus constructed, we find the first non-trivial vector of coefficients to be
\begin{eqnarray}
&&\{-1334,-860,1081,908,-504,95,-54,-359,28,0,188,272,225,1013,-280,0,-985,332,168,0,-1848,0,696,
\nonumber\\&&\ \  160,-712,-756,466,-240,-160,-93,408,-58,785,348,-408,-2\}\,,
\label{eq:one}
\end{eqnarray}
where the last entry is the normalization factor.
Eq.~\eqref{eq:one} is far to complicated to be true.

As is well known from previous results, the coefficients of the first few basis elements in Eq.~\eqref{BasisCF3nf}, namely, the lower powers of $\eta$ or $\nu$ and the simplest binomial harmonic sums like $\BS_{1}$, are relatively ``complicated'' numbers upon factoring out $2^i3^j$.
Therefore, the chances of \texttt{fplll} to succeed may be increased by eliminating one of them using one out of the nine available equations.
Eliminating the first basis element, $1$, using the first equation, for $N=4$, we obtain
\begin{eqnarray}
&&\{-889,199,-1816,1275,37,-419,-107,-56,-280,-275,-210,107,-194,-366,0,-568,181,-336,96,6,-24,\nonumber\\&&\ \
  -348,108,7,709,-204,262,-132,48,-23,178,-158,1030,26,4\}\,,
\label{eq:two}
\end{eqnarray}
which does not qualify.
Eliminating, one after another, also $\eta,\nu,\eta^2$ using, in this order, the equations for $N=6,8,10$, we get
\begin{eqnarray}
  &&\{-731,908,-264,-132,601,100,28,24,-169,106,-473,-638,176,0,-449,154,168,0,-246,0,954,80,255,
  \nonumber\\&&\ \
  177,80,-84,-80,45,177,-200,-320,-354,120,-2\}\,,\nonumber\\
&&\{454,30,-205,88,-364,14,84,178,-129,-11,-493,-372,0,219,-73,84,48,176,-12,364,-136,-13,48,\nonumber\\&&\ \
  -306,-232,124,107,240,126,-104,-288,180,-1\}\,,\nonumber\\
 &&\{85,-205,-254,95,14,136,19,-144,180,135,-153,0,123,-41,84,-36,-108,9,51,-88,-113,21,91,\nonumber\\&&\ \
  -177,97,107,-75,1,-4,54,-92,-1\}\,,
\label{eq:three}
\end{eqnarray}
none of which qualifies either.

However, if we trade $\nu$ by $\BS_{1}$, we find
\begin{eqnarray}
&&\{-27,48,219,-84,14,0,354,-24,30,132,-56,0,-66,24,84,0,192,0,216,0,-12,0,0,-60,0,-12,0,0,\nonumber\\&&\ \
14,0,0,-1\}\,,
\label{eq:four}
\end{eqnarray}
which strikes as the solution.
In fact, the vector in Eq.~\eqref{eq:four} distinctly differs from those in Eqs.~\eqref{eq:one}--\eqref{eq:three}, by carrying numerous zeros alongside integers that can be factored as $2^i3^j$.
This qualitative breakthrough serves as the primary argument in favor of Eq.~\eqref{eq:four}.

In the case at hand, we still have the equation for $N=22$ \cite{Falcioni:2025hfz} and the Chinese remainder theorem \cite{china}, fixing the coefficients of $\eta^{4,3,2}$ and $\nu^4$ \cite{Moch:2014sna}, for cross checks, which all work out all right.
In other cases, when all available information is just enough to produce a solution with such distinctive properties as in Eq.~\eqref{eq:four}, we are still on the safe side.

\boldmath
\subsection{All-$N$ results for $\gamma_{gg,\zeta_k}^{(3)}$ with $k=3,4,5$}
\label{supplement:2}
\unboldmath

For the reader's convenience, we list here the complete all-$N$ expressions for the transcendental contributions $\gamma_{gg,\zeta_k}^{(3)}$ with $k=3,4,5$ in
Eq.~\eqref{eq:zeta3}:


\begin{eqnarray}
&&\gamma_{gg,\zeta_3}^{(3)}=





\frct{32}{3}\*C_F^3\*n_f\*
\Big[
24\*\D1
-24\*\D2
+231\*\D1^2
+294\*\D1^3
+144\*\D1^4
-\frct{897}{4}\*\eta
-16\*\nu
-177\*\eta^2
-138\*\eta^3
-114\*\eta^4
+6\*(
6\*\D1
\nonumber\\
&&{}
-6\*\D2
-2\*\D1^2
+8\*\D1^3
-71\*\eta
+20\*\nu
-78\*\eta^2
-56\*\eta^3
)\*\HS_{1}
-4\*(15\*\eta
-4\*\nu
+18\*\eta^2
)\*\HS_{2}
+4\*(33\*\eta
-16\*\nu
+90\*\eta^2
)\*\HS_{-2}
+12\*(13\*\eta
\nonumber\\
&&{}-4\*\nu
+14\*\eta^2
)\*\HS_{1,1}
-\frct{2}{3}\*\delta_{N,2}
\Big]

+\frct{32}{9}\*C_A\*C_F^2\*n_f\*\Big\{
-\frct{11}{2}
+\frct{1161}{2}\*(\D1-\D2)
-2817\*\D1^2
-1476\*\D1^3
-1836\*\D1^4
+148\*\D2^2
\nonumber\\
&&{}
+120\*\D2^3
+176\*\nu^2
+\frct{6879}{2}\*\eta
-\frct{1597}{3}\*\nu
+\frct{3771}{2}\*\eta^2
+2376\*\eta^3
+1224\*\eta^4
+3\*\Big[\frct{37}{2}
-30\*(\D1-\D2)
+123\*\D1^2
+12\*\D2^2
+132\*\D1^3
\nonumber\\
&&{}
+1065\*\eta
-366\*\nu
+1656\*\eta^2
-8\*\nu^2
+930\*\eta^3
\Big]\*\HS_{1}
+3\*(231\*\eta
-68\*\nu
+330\*\eta^2
)\*\HS_{2}
-24\*(24\*\eta
-13\*\nu
+81\*\eta^2
)\*\HS_{-2}
-48\*(30\*\eta
-10\*\nu
\nonumber\\
&&{}
+39\*\eta^2
)\*\HS_{1,1}
+\frct{7}{3}\*\delta_{N,2}
\Big\}

+\frct{32}{3}\*C_A^2\*C_F\*n_f\*\Big\{\frct{2}{3}
-\frct{471}{2}\*(\D1-\D2)
+\frct{2019}{2}\*\D1^2
-\frct{514}{3}\*\D2^2
+201\*\D1^3
-136\*\D2^3
+684\*\D1^4
\nonumber\\
&&{}
-\frct{2530}{3}\*\eta
+\frct{1303}{18}\*\nu
-\frct{1325}{3}\*\eta^2
-\frct{256}{3}\*\nu^2
-892\*\eta^3
-40\*\nu^3
-420\*\eta^4
+\Big[\frct{241}{6}
-6\*(\D1-\D2)
-135\*\D1^2
-84\*\D2^2
-324\*\D1^3
\nonumber\\
&&{}
-552\*\eta
+\frct{773}{3}\*\nu
-1470\*\eta^2
+112\*\nu^2
-726\*\eta^3
\Big]\*\HS_{1}
-(
309\*\eta
-124\*\nu
+426\*\eta^2
)\*\HS_{2}
+2\*(3
-3\*\eta
-8\*\nu
+210\*\eta^2
)\*\HS_{-2}
+8\*(60\*\eta
\nonumber\\
&&{}
-23\*\nu
+90\*\eta^2
)\*\HS_{1,1}
+8\*\HS_{3}
+8\*\HS_{-3}
-24\*\HS_{1,2}
-24\*\HS_{1,-2}
-24\*\HS_{2,1}
+8\*\HS_{-2,1}
-\frct{2}{9}\*\delta_{N,2}
\Big\}

+\frct{32}{3}\*C_A^3\*n_f\*\Big[
\frct{2197}{360}
-294\*\D1^2
\nonumber\\
&&{}
+122\*\D2^2
-24\*\D1^3
+96\*\D2^3
-216\*\D1^4
+\frct{683}{12}\*\eta
+\frct{6734}{135}\*\nu
+\frct{2431}{12}\*\eta^2
+\frct{194}{3}\*\nu^2
+377\*\eta^3
+20\*\nu^3
+144\*\eta^4
+\frct{1}{18}\*(
-1173
+432\*\D1^2
\nonumber\\
&&{}
+1296\*\D2^2
+2592\*\D1^3
+4989\*\eta
-1696\*\nu
+8913\*\eta^2
-1440\*\nu^2
+4194\*\eta^3
)\*\HS_{1}
+\frct{5}{3}\*(99\*\eta
-40\*\nu
+54\*\eta^2
)\*\HS_{2}
-\frct{2}{3}\*(
315\*\eta
-92\*\nu
\nonumber\\
&&{}
+162\*\eta^2
)\*\HS_{1,1}
+\frct{1}{6}\*\Big(
\frct{128}{15}
+\frct{2436}{5}\*\eta
-116\*\nu
-1056\*\eta^2
\Big)\*\HS_{-2}
-\frct{1}{6}\*(
59
-18\*\eta)\*\HS_{3}
+\frct{1}{3}\*(31
-18\*\eta)\*\HS_{-3}
+24\*\HS_{1,2}
+\frct{46}{3}\*\HS_{1,-2}
\nonumber\\
&&{}
+24\*\HS_{2,1}
+12\*(
-3
+\eta)\*\HS_{-2,1}
+\frct{1}{2}\*N\*(N+1)\*\Big(
\frct{4}{45}\*\HS_{-2}
+\HS_{3}
-2\*\HS_{-3}
+4\*\HS_{-2,1}\Big)\Big]

+\frct{64}{3}\*C_A^4\*\Big[
-\frct{2621}{720}
-\frct{36013}{144}\*\eta
+\frct{101027}{1080}\*\nu
\nonumber\\
&&{}
-\frct{1133}{6}\*\eta^2
-\frct{207}{4}\*\nu^2
-97\*\eta^3
+19\*\nu^3
-12\*\eta^4
-24\*\nu^4
+\frct{1}{36}\*(437
-21021\*\eta
+5915\*\nu
-12174\*\eta^2
-2664\*\nu^2
-4320\*\eta^3
\nonumber\\
&&{}
+1728\*\nu^3)\*\HS_{1}
+\frct{1}{3}\*(5
-51\*\eta
+33\*\nu)\*\HS_{3}
+6\*\HS_{4}
+\frct{4}{3}\*(45\*\eta
+45\*\eta^2
-25\*\nu
+12\*\nu^2)\*(\HS_{2}-2\*\HS_{1,1})
-\frct{1}{18}\*\Big(
\frct{452}{5}
-\frct{1023}{5}\*\eta
+435\*\nu
\nonumber\\
&&{}
-612\*\eta^2
+72\*\nu^2
\Big)\*\HS_{-2}
-\frct{1}{3}\*(
131
-141\*\eta
+105\*\nu)\*\HS_{-3}
-18\*\HS_{-4}
-11\*\HS_{1,3}
+\frct{1}{6}\*(143
-96\*\eta
+96\*\nu)\*\HS_{1,-2}
+(127
-156\*\eta
\nonumber\\
&&{}
+108\*\nu)\*\frct{1}{2}\*\HS_{-2,1}
+35\*\HS_{1,-3}
+8\*\HS_{2,-2}
-5\*\HS_{3,1}
+4\*\HS_{-2,-2}
+16\*\HS_{-2,2}
+43\*\HS_{-3,1}
-16\*\HS_{1,1,-2}
-54\*\HS_{1,-2,1}
-32\*\HS_{-2,1,1}
\nonumber\\
&&{}
-N\*(N+1)\*\Big(
\frct{1}{90}\*\HS_{-2}
+\frct{1}{4}\*\HS_{3}
-\frct{1}{2}\*\HS_{-3}
+\HS_{-2,1}
\Big)
\Big]

+\frct{64}{3}\*\frct{d_{AA}^{(4)}}{n_a}\*
  \Big[
 \frct{331}{15}
+ \frct{429}{2}\*\eta
- \frct{3824}{45}\*\nu
+ 180\*\eta^2
+ 2\*\HS_{1} \*\Big(-5 + 1352\*\eta - \frct{1328}{3}\*\nu
\nonumber\\
&&{}
+ 722\*\eta^2
+ 192\*\nu^2
+ 144\*\eta^3
\Big)
- 4\*\HS_{-2}\*\Big(\frct{1126}{15}+\frct{1127}{5}\*\eta
- 19\*\nu + 48\*\nu^2 - 12\*\eta^2 \Big)
- 16\*(2\*\HS_{1,1}-\HS_{2})\*(63\*\eta - 25\*\nu + 54\*\eta^2 +12\*\nu^2)
\nonumber\\
&&{}
+4\*(31-102\*\eta+30\*\nu)\*\HS_{-3}
-4\*(43-66\*\eta+30\*\nu)\*\HS_{3}
-36\*(13-28\*\eta+12\*\nu)\*\HS_{-2,1}
+4\*(55-48\*\eta+48\*\nu)\*\HS_{1,-2}
-24\*\HS_{-3,1}
\nonumber\\
&&{}
-96\*\HS_{-2,-2}
+192\*\HS_{-2,2}
-120\*\HS_{1,-3}
+120\*\HS_{1,3}
+96\*\HS_{2,-2}
-24\*\HS_{3,1}
-384\*\HS_{-2,1,1}
+432\*\HS_{1,-2,1}
-192\*\HS_{1,1,-2}
\nonumber\\
&&{}
-2\* N\*(N+1)\*
   \Big(
\frct{1}{15}\* \HS_{-2}
+ 6\* \HS_{-3}
- 3\* \HS_{3}
- 12\* \HS_{-2, 1}
\Big)
\Big]


+\frct{128}{3}\*\frct{d_{RA}^{(4)}}{n_a}\*n_f\*\Big[
\frct{103}{15}
+\frct{2908}{45}\*\nu
-291\*\eta
-387\*\eta^2
+2\*\Big(
5
+ 14\*\eta
+\frct{52}{3}\*\nu
\nonumber\\
&&{}
- 46\*\eta ^2 - 84\*\eta ^3\Big)\*\HS_{1}
+2\*\Big(\frct{514}{15}+\frct{618}{5}\*\eta
- 4\*\nu
- 132\*\eta^2\Big)\*\HS_{-2}
+2\*\big(9\*\eta - 16\*\nu + 198\*\eta^2\big)\*(2\*\HS_{1,1}-\HS_{2})
- 4\* (13 - 18\* \eta)\* \HS_{-3}
\nonumber\\
&&{}
+ 2\* (23 - 18\* \eta)\* \HS_{3}
+ 144\* (1 - \eta)\* \HS_{-2, 1}
- 40\* \HS_{1, -2}
+(N+1)\*N\*\Big(
\frct{4}{15}\* \HS_{-2}
+ 12\* \HS_{-3}
- 6\* \HS_{3}
- 24\* \HS_{-2, 1}
\Big)
\Big]
+\frct{16}{27}\*C_A^2\*n_f^2\*\Big\{
-\frct{13}{10}
\nonumber\\
&&{}
+\frct{4773}{4}\*\eta
-\frct{7154}{15}\*\nu
+\frct{2577}{2}\*\eta^2
-84\*\nu^2
+324\*\eta^3
-\frct{2}{5}\*
[
N\*(N+1)
+17
+42\*\eta
]\*\HS_{-2}
+
(120
-753\*\eta
+230\*\nu
-492\*\eta^2
+144\*\eta^3
)\*\HS_{1}
\nonumber\\
&&{}
-36\*\eta\*
(
1
-4\*\eta)\*\HS_{2}
+72\*\eta\*
(
1
-4\*\eta)\*\HS_{1,1}\Big\}

+\frct{256}{9}\*\frct{d_{RR}^{(4)}}{n_a}\*n_f^2\*\Big\{
-\frct{29}{5}
+\frct{297}{2}\*\eta
-\frct{299}{15}\*\nu
+135\*\eta^2
-\frct{2}{5}\*
[17
+42\*\eta
+N\*(N+1) ] \*\HS_{-2}
\nonumber\\
&&{}
-2\*
(
57\*\eta
-\nu
-78\*\eta^2
-72\*\eta^3
)\*\HS_{1}
-36\*\eta\*
(
1
-4\*\eta)\*\HS_{2}
+72\*\eta\*
(
1
-4\*\eta)\*\HS_{1,1}\Big\}

+\frct{32}{9}\*C_A\*C_F\*n_f^2\*\Big[
-5
-21\*(\D1-\D2)
-\frct{799}{2}\*\eta
\nonumber\\
&&{}
+143\*\nu
+15\*\D1^2
-588\*\eta^2
+24\*\nu^2
-12\*\D1^3
-150\*\eta^3
-4\*
(
5
-54\*\eta
+19\*\nu
-60\*\eta^2
)\*\HS_{1}
+\frct{1}{3}\*\delta_{N,2}
\Big]

+\frct{32}{9}\*C_F^2\*n_f^2\*\Big[
48\*(\D1-\D2)
\nonumber\\
&&{}
+\frct{11}{2}
+231\*\eta
-\frct{214}{3}\*\nu
-24\*\D1^2
+395\*\eta^2
-8\*\nu^2
+48\*\D1^3
+90\*\eta^3
-2\*
(63\*\eta
-20\*\nu
+66\*\eta^2
)\*\HS_{1}
\Big]

+
\frct{64}{27}\*C_A\*n_f^3\*
(\eta
-\nu
+\HS_{1})

\nonumber\\
&&{}
+\frct{32}{9}\*C_F\*n_f^3\*
(9\*\eta
-4\*\nu
+6\*\eta^2)


\,,
\\



&&\gamma_{gg,\zeta_4}^{(3)}=






176\*C_A^3\*n_f\*
(4\*\eta
-2\*\nu
+3\*\eta^2
+\HS_{1})

-8\*C_A^2\*C_F\*n_f\*
[25\*\eta
-12\*\nu
+60\*\eta^2
-16\*\nu^2
-24\*\eta^3
+2\*
(11
-18\*\eta
+8\*\nu
-12\*\eta^2
)\*\HS_{1}]
\nonumber\\
&&{}

-16\*C_A\*C_F^2\*n_f\*
[12\*\eta
-8\*\nu
-15\*\eta^2
+8\*\nu^2
+6\*\eta^3
+4\*
(9\*\eta
-4\*\nu
+6\*\eta^2
)\*\HS_{1}]

-8\*C_F^3\*n_f\*
[39\*\eta
-16\*\nu
+36\*\eta^2
+12\*\eta^3
-4\*
(9\*\eta
-4\*\nu
\nonumber\\
&&{}
+6\*\eta^2
)\*\HS_{1}]

-32\*C_A^2\*n_f^2\*
(4\*\eta
-2\*\nu
+3\*\eta^2
+\HS_{1})

+\frct{16}{3}\*C_A\*C_F\*n_f^2\*
(15\*\eta
-4\*\nu
+24\*\eta^2
+6\*\HS_{1})

+\frct{16}{3}\*C_F^2\*n_f^2\*
(9\*\eta
-8\*\nu
-6\*\eta^2
   )


\,,
\\




&&\gamma_{gg,\zeta_5}^{(3)}=






\frct{160}{9}\*\Big(C_A^4+12\*\frct{d_{AA}^{(4)}}{n_a}\Big)\*
\Big[
-\frct{1}{60}\*N\*(N+1)
-\frct{751}{30}
-21\*\nu
+\frct{94}{5}\*\eta
-12\*\nu^2
+36\*\eta^2
+(33
+24\*\nu
-24\*\eta)\*\HS_{1}
+12\*\HS_{2}
\nonumber\\ &&{}
-24\*\HS_{1,1}
\Big]
+\frct{80}{9}\*C_A^3\*n_f\*
\Big[
\frct{1}{15}\*N\*(N+1)
+\frct{17}{15}
+16\*\nu
+\frct{149}{5}\*\eta
-456\*\eta^2
+6\*\HS_{1}
\Big]

+\frct{640}{3}\*\frct{d_{RA}^{(4)}}{n_a}\*n_f\*\Big[
\frct{287}{15}
+\frct{1}{15}\*N\*(N+1)
+4\*\nu
\nonumber\\ &&{}
-\frct{121}{5}\*\eta
+30\*\eta^2
-12\*\HS_{1}
\Big]

+80\*C_A^2\*C_F\*n_f\*\Big(2
-\frct{68}{3}\*\nu
+41\*\eta
+184\*\eta^2
+2\*\HS_{1}
\Big)

+80\*C_A\*C_F^2\*n_f\*\big(32\*\nu
-81\*\eta
-262\*\eta^2
-4\*\HS_{1}
\big)
\nonumber\\ &&{}

+160\*C_F^3\*n_f\*\Big(
-\frct{16}{3}\*\nu
+17\*\eta
+66\*\eta^2
\Big)



+\frct{8}{9}\*\Big(C_A^2+48\*\frct{d_{RR}^{(4)}}{n_a}\Big)\*n_f^2\*\Big[
-\frct{1}{3}\*N\*(N+1)
-\frct{17}{3}
+46\*\eta
-240\*\eta^2
\Big]


\,,
\end{eqnarray}
where the notation is as in the main text.
\end{widetext}

\end{document}
===== END FILE: 2604.25833v2.tex =====
